%%
%% The original source files were:
%%
%% samples.dtx  (with options: `all,proceedings,bibtex,manuscript')
%% 
%% IMPORTANT NOTICE:
%% 
%% For the copyright see the source file.
%% 
%% Any modified versions of this file must be renamed
%% with new filenames distinct from sample-manuscript.tex.
%% 
%% For distribution of the original source see the terms
%% for copying and modification in the file samples.dtx.
%% 
%% This generated file may be distributed as long as the
%% original source files, as listed above, are part of the
%% same distribution. (The sources need not necessarily be
%% in the same archive or directory.)
%%
%%
%% Commands for TeXCount
%TC:macro \cite [option:text,text]
%TC:macro \citep [option:text,text]
%TC:macro \citet [option:text,text]
%TC:envir table 0 1
%TC:envir table* 0 1
%TC:envir tabular [ignore] word
%TC:envir displaymath 0 word
%TC:envir math 0 word
%TC:envir comment 0 0
%%
%% The first command in your LaTeX source must be the \documentclass
%% command.
%%
%% For submission and review of your manuscript please change the
%% command to \documentclass[manuscript, screen, review]{acmart}.
%%
%% When submitting camera ready or to TAPS, please change the command
%% to \documentclass[sigconf]{acmart} or whichever template is required
%% for your publication.
%%
%%
\documentclass[manuscript,screen,review]{acmart}


\usepackage{algorithm}
% Use algorithmicx + algpseudocode: provides \Require, \Ensure, \State, and \Comment
\usepackage{algorithmicx}
\usepackage{algpseudocode}
% Provide compatibility shims for uppercase macros used in the document
% Some authors use \REQUIRE/\ENSURE/\UPON/\ENDUPON (uppercase) — map them to algpseudocode
\providecommand{\REQUIRE}[1]{\Require{#1}}
\providecommand{\ENSURE}[1]{\Ensure{#1}}
% \UPON{<cond>} will render as a bold 'Upon <cond>' state line inside algorithmic blocks
\providecommand{\UPON}[1]{\State\textbf{Upon}~#1}
\providecommand{\ENDUPON}{}
% Provide uppercase aliases commonly used in algorithm blocks
\providecommand{\STATE}{\State}
\providecommand{\COMMENT}[1]{\Comment{#1}}
% Map common uppercase algorithmic keywords to algorithmicx/algpseudocode commands
\providecommand{\IF}{\If}
\providecommand{\ELSE}{\Else}
\providecommand{\ELSIF}{\ElsIf}
\providecommand{\ENDIF}{\EndIf}
\providecommand{\FOR}{\For}
\providecommand{\ENDFOR}{\EndFor}
\providecommand{\WHILE}{\While}
\providecommand{\ENDWHILE}{\EndWhile}
\providecommand{\RETURN}{\State\textbf{return}~}
\usepackage{tikz}
% Load common tikz libraries used in figures (positioning provides the "of" syntax)
\usetikzlibrary{positioning,arrows.meta,calc,shapes,shadows}
% color support for \textcolor used inside algorithms
\usepackage{xcolor}
% Math packages required for \text and bold symbols inside math
\usepackage{amsmath}
\usepackage{bm}
\usepackage{adjustbox}
% Define common math operators used in the document
\DeclareMathOperator*{\argmin}{arg\,min}
\DeclareMathOperator*{\argmax}{arg\,max}
\usepackage{multirow}
% Enable list customization (leftmargin, itemsep, etc.) used in the document
\usepackage{enumitem}
\newlength{\tblwidth}
\setlength{\tblwidth}{\linewidth}
% Ensure header height is sufficient for fancyhdr
\setlength{\headheight}{20.48303pt}
% theorem environments used in the document
\usepackage{amsthm}
\newtheorem{theorem}{Theorem}
\newtheorem{definition}{Definition}
\newtheorem{assumption}{Assumption}
\newtheorem{problem}{Problem}
\newtheorem{lemma}{Lemma}
\newtheorem{corollary}{Corollary}
\newtheorem{remark}{Remark}

\newenvironment{tablenotes}{%
    \par\footnotesize
    \begin{list}{}{% 
        \setlength{\leftmargin}{1.5em}%
        \setlength{\labelsep}{0.5em}%
        \setlength{\itemsep}{0pt}%
        \setlength{\parsep}{0pt}%
    }%
}{%
    \end{list}%
}

%%
%% \BibTeX command to typeset BibTeX logo in the docs
\AtBeginDocument{%
  \providecommand\BibTeX{{%
    Bib\TeX}}}

%% Rights management information.  This information is sent to you
%% when you complete the rights form.  These commands have SAMPLE
%% values in them; it is your responsibility as an author to replace
%% the commands and values with those provided to you when you
%% complete the rights form.
\setcopyright{acmlicensed}
\copyrightyear{2025}
\acmYear{2025}
\acmDOI{XXXXXXX.XXXXXXX}

%% These commands are for a PROCEEDINGS abstract or paper.
%%  \acmConference[Conference acronym 'XX]{Make sure to enter the correct
%% 	conference title from your rights confirmation email}{June 03--05,
%% 	2018}{Woodstock, NY}

%%
%%  Uncomment \acmBooktitle if the title of the proceedings is different
%%  from ``Proceedings of ...''!
%%
%%\acmBooktitle{Woodstock '18: ACM Symposium on Neural Gaze Detection,
%%  June 03--05, 2018, Woodstock, NY}
%%\acmISBN{978-1-4503-XXXX-X/2018/06}


%%
%% Submission ID.
%% Use this when submitting an article to a sponsored event. You'll
%% receive a unique submission ID from the organizers
%% of the event, and this ID should be used as the parameter to this command.
%%\acmSubmissionID{123-A56-BU3}

%%
%% For managing citations, it is recommended to use bibliography
%% files in BibTeX format.
%%
%% You can then either use BibTeX with the ACM-Reference-Format style,
%% or BibLaTeX with the acmnumeric or acmauthoryear sytles, that include
%% support for advanced citation of software artefact from the
%% biblatex-software package, also separately available on CTAN.
%%
%% Look at the sample-*-biblatex.tex files for templates showcasing
%% the biblatex styles.
%%

%%
%% The majority of ACM publications use numbered citations and
%% references.  The command \citestyle{authoryear} switches to the
%% "author year" style.
%%
%% If you are preparing content for an event
%% sponsored by ACM SIGGRAPH, you must use the "author year" style of
%% citations and references.
%% Uncommenting
%% the next command will enable that style.
%%\citestyle{acmauthoryear}


%%
%% end of the preamble, start of the body of the document source.
\begin{document}

%%
%% The "title" command has an optional parameter,
%% allowing the author to define a "short title" to be used in page headers.
\title{Kinematic Spatiotemporal Constraint-Based Detection and Adaptive Isolation of Byzantine False Messages in Industrial Multi-Agent Systems}

%%
%% The "author" command and its associated commands are used to define
%% the authors and their affiliations.
%% Of note is the shared affiliation of the first two authors, and the
%% "authornote" and "authornotemark" commands
%% used to denote shared contribution to the research.
\author{Qi Yu}
%%\authornote{.}
%%\authornotemark[1]
\email{230240012@seu.edu.cn}
%%\orcid{1234-5678-9012}
\affiliation{%
	\institution{School of Cyber Science and Engineering, Southeast University}
	\city{Nanjing}
	\state{Jiangsu}
	\country{PR China}
}


\author{Wei Li}
%%\orcid{1234-5678-9012}
\email{xchlw@seu.edu.cn}
\authornote{corresponding author}
%%\authornotemark[1]
\affiliation{%
	\institution{School of Computer Science and Engineering, Southeast University}
	\city{Nanjing}
	\state{Jiangsu}
	\country{PR China}
}


%%
%% By default, the full list of authors will be used in the page
%% headers. Often, this list is too long, and will overlap
%% other information printed in the page headers. This command allows
%% the author to define a more concise list
%% of authors' names for this purpose.
%%\renewcommand{\shortauthors}{Qi et al.}

%%
%% The abstract is a short summary of the work to be presented in the
%% article.
\begin{abstract}
In large-scale industrial multi-agent systems, Byzantine agents can inject physically implausible messages---reporting false positions, forged timestamps, or kinematically infeasible velocities---to disrupt coordination and endanger physical safety. Existing Byzantine Fault Tolerance (BFT) protocols treat messages as abstract data and rely solely on cryptographic authentication, failing to exploit the rich physical constraints inherent to industrial cyber-physical systems (CPS). This paper proposes a \textit{Kinematic Spatiotemporal Constraint-Based Detection and Adaptive Isolation} (KSCDI) framework that leverages three categories of physical invariants---temporal bounds (clock skew $\leq \delta_t$), kinematic feasibility ($|\Delta\mathbf{p}| \leq v_{\max} \cdot \Delta t$), and task-membership consistency---to detect and reject Byzantine false messages before they propagate through the system. Detected Byzantine nodes are progressively isolated via adaptive trust weights that decay under repeated violations and recover under sustained honest behavior, without requiring explicit membership reconfiguration. We formally characterize the detection coverage and isolation convergence under both random and task-concentrated Byzantine fault distributions, and clarify the assumptions on sensor noise bounds and failure detection completeness. Experimental validation on a Raspberry Pi hardware testbed (6 agent nodes) demonstrates 100\% detection of timestamp forgery and kinematic violations, 98\% detection of position spoofing (with one false negative attributable to GPS noise of $\pm$3\,m), and Byzantine node weight isolation within 4--7 rounds. Emulation-based evaluation (CORE, 10--100 nodes) further confirms 90\% early message filtering under 30\% Byzantine load and 168\,ms consensus latency, well within the 500\,ms real-time threshold for industrial applications. These results demonstrate that kinematic spatiotemporal constraints provide a lightweight, cryptography-free detection layer that significantly reduces Byzantine influence in resource-constrained industrial multi-agent systems.
\end{abstract}

\keywords{Byzantine fault detection, Kinematic constraint validation, Spatiotemporal message filtering, Adaptive trust isolation, Industrial cyber-physical systems, Industrial multi-agent systems, False message rejection, Byzantine fault tolerance}

\maketitle


%=============================================================================
\section{Introduction}
\label{sec:introduction}

Byzantine Fault Tolerance (BFT) ensures reliable consensus in distributed systems where nodes may exhibit arbitrary malicious behavior~\cite{wu2024half, liu2024improved}. In industrial cyber-physical systems (CPS)---such as smart factories and autonomous warehouses---Byzantine agents can inject physically implausible messages: reporting false positions, forging timestamps, or claiming kinematically infeasible velocities to disrupt coordination and endanger physical safety~\cite{bhamare2020cybersecurity}. Existing BFT protocols treat messages as abstract data and rely solely on cryptographic authentication, failing to exploit the rich physical constraints inherent to industrial environments. This paper addresses the question: \textit{can kinematic and spatiotemporal constraints be used as a lightweight, cryptography-free detection layer to filter Byzantine false messages before they propagate through the consensus protocol?}

\subsection{Motivation and Key Insight}

Modern smart factories and autonomous warehouses deploy large fleets of heterogeneous agents that collaborate to execute complex production workflows~\cite{vlachos2024smart}. Consider a typical material handling scenario in a smart factory with agents (Figure~\ref{fig:motivation}):
\begin{itemize}
        \item \textbf{Automated Guided Vehicles (AGVs)} transport raw materials and semi-finished products between workstations, navigating autonomously using LiDAR and visual odometry.
        \item \textbf{Robot Arms (RAs)} perform assembly, welding, and quality inspection at fixed workstations, equipped with force sensors and vision systems.
        \item \textbf{Scheduling System (S)} coordinates task allocation, monitors production progress, and optimizes resource utilization.
\end{itemize}

\begin{figure*}[!t]
        \centering
        \includegraphics[width=.85\textwidth]{fig/motivation_example.png}
        \Description{Smart factory scenario illustrating Byzantine false message threats, where a compromised AGV reports a kinematically infeasible position and the spatiotemporal validation layer rejects it before consensus.}
        \caption{Smart factory scenario illustrating Byzantine false message threats. A compromised AGV reports a kinematically infeasible position (teleportation beyond $v_{\max} \cdot \Delta t$). KSCDI's spatiotemporal validation layer detects and rejects this message before it enters consensus, preventing physical collision.}
        \label{fig:motivation}
\end{figure*}

A typical task involves \textit{spatially localized collaboration}: AGV$_1$ transports a component to Workstation~A, where Robot Arm RA$_3$ retrieves it for assembly. This handoff requires precise coordination: AGV$_1$ confirms arrival at the correct position, RA$_3$ verifies component identity, and both agree on the handoff timestamp. The decision must update the production database reliably. Critically, this coordination involves only \textit{three agents} (AGV$_1$, RA$_3$, and the Scheduling System~S), yet the decision must be \textit{fault-tolerant} to ensure production correctness and safety.

\subsubsection{Byzantine Threats in Industrial Environments}

Industrial multi-agent systems face diverse Byzantine threats that go beyond simple crash failures~\cite{bhamare2020cybersecurity}:

\textbf{Malicious Agents.} Compromised or counterfeit devices may send conflicting state reports to disrupt production. For example, a hacked AGV could report false positions, or a malicious robot arm could claim successful assembly when parts are defective.

\textbf{Network Attacks.} Man-in-the-middle attacks can alter messages between agents, replay old messages to cause duplicate operations, or selectively drop messages to delay critical tasks.

The consequences of Byzantine failures in industrial settings are severe: incorrect material routing causes production delays~\cite{margret2024blockchain}, undetected defective products lead to safety recalls, and coordinated attacks on AGV fleets can cause physical collisions endangering human workers.

Unlike abstract distributed systems where Byzantine behavior is arbitrary, industrial agents operate in the \textit{physical world} with measurable properties: positions are constrained by kinematics (an AGV cannot teleport), velocities are bounded by motor capabilities, and timestamps follow causal ordering~\cite{li2022byzantine}. These physical constraints provide strong, cryptography-free validation signals---but existing BFT protocols fail to exploit them.

\textbf{Remark on Fault Distribution.} We consider both uniformly distributed Byzantine faults and task-concentrated faults, where an adversary compromises multiple agents assigned to the same task (e.g., targeting a specific workstation area). The latter is arguably more realistic in industrial deployments and is explicitly evaluated in Section~\ref{sec:evaluation}.

\subsubsection{Why Byzantine Fault-Tolerant Consensus Helps}

A critical subtlety in industrial BFT is clarifying \textit{what agents agree on} and \textit{how agreement prevents physical harm}. Consider the following concrete scenario:

\textbf{Scenario.} AGV$_1$ is compromised and broadcasts a false arrival position $\hat{p} = (x', y')$ that deviates 3\,m from its true location~$(x, y)$. If RA$_3$ blindly trusts this single report and initiates its retrieval motion toward~$(x', y')$, a physical collision with AGV$_1$ (still at $(x, y)$) becomes inevitable---even if RA$_3$ acts in good faith.

\textbf{What STLC agents agree on.} Rather than trusting any single agent's self-report, STLC runs Byzantine consensus within the task-coupled domain $C(\tau_j) = \{\text{AGV}_1, \text{RA}_3, S\}$. The \textit{decided value} is a tuple of verified task execution parameters:
\[
    v^* = \langle p^*_{\text{AGV}_1},\; t^*_{\text{handoff}},\; \text{id}^*_{\text{component}} \rangle
\]
where $p^*_{\text{AGV}_1}$ is the \emph{consensus-agreed} arrival position of AGV$_1$, $t^*_{\text{handoff}}$ is the agreed handoff timestamp, and $\text{id}^*_{\text{component}}$ is the agreed component identifier. This agreed value is derived from cross-validated reports of all honest agents in the domain, not from any single node's claim.

\textbf{How physical constraint detection prevents harm.} Before entering consensus, STLC's spatiotemporal validation layer checks whether AGV$_1$'s reported position $\hat{p}$ satisfies the kinematic feasibility constraint:
\[
    |\hat{p} - p_{\text{prev}}| \leq v_{\max} \cdot \Delta t
\]
The false report $(x', y')$ violates this constraint and is rejected immediately, without cryptographic operations. The remaining honest agents (RA$_3$ and~$S$) then reach consensus on the true position $(x, y)$, which RA$_3$ uses as the authoritative handoff location. Even if spatiotemporal filtering is bypassed by a sophisticated attacker, the BFT quorum requirement ensures that a single Byzantine agent cannot unilaterally determine $v^*$: at least $\lfloor 2k/3 \rfloor + 1$ honest agents must agree on the decided value.

\textbf{Detection completeness assumption.} We assume that correct agents report positions and timestamps within bounded sensor noise (GPS error $\leq \pm 3$\,m, clock drift $\leq 100$\,ms). Under this assumption, the spatiotemporal filter achieves high but not perfect detection: one false negative was observed in 50 trials due to a GPS noise spike at the boundary of the tolerance threshold (Section~\ref{sec:evaluation}). Byzantine agents that remain within sensor noise bounds are handled by the BFT quorum layer.

\subsubsection{Key Insight: Exploiting Physical Constraints for Byzantine Detection}

Industrial agents operate under physical laws that provide strong validation signals unavailable in abstract distributed systems: \textbf{(1) Kinematic constraints}---an AGV cannot teleport; its reported position must be reachable from its last confirmed location within the elapsed time; \textbf{(2) Temporal causality}---timestamps must respect clock skew bounds and causal ordering; \textbf{(3) Task-membership consistency}---only agents assigned to a task can legitimately participate in its consensus.

These three invariants enable a lightweight pre-consensus filtering layer that rejects 65--90\% of Byzantine messages before they consume consensus resources, without requiring cryptographic operations. Nodes that repeatedly violate these constraints are progressively isolated via adaptive trust weights that decay under violations and recover under sustained honest behavior---without explicit membership reconfiguration.

\subsection{Contributions}
We propose \textbf{STLC} (Spatiotemporal Localized Consensus with Adaptive Trust) with the following contributions:

\textbf{(1) Kinematic Spatiotemporal Constraint Detection.} We formalize three categories of physical invariants---temporal bounds, kinematic feasibility, and task-membership consistency---as a pre-consensus Byzantine detection layer. We characterize detection coverage and the conditions under which Byzantine messages evade detection (sensor noise boundary cases), and analyze detection completeness under both random and task-concentrated fault distributions.

\textbf{(2) Adaptive Trust Isolation.} We design a trust weight mechanism that progressively isolates nodes exhibiting repeated constraint violations via weighted quorum voting. Byzantine nodes reach isolation weight $w_{\min}$ within 4--7 rounds; honest nodes with transient faults recover within 10 rounds. Isolation occurs without explicit membership changes, avoiding the coordination overhead of view-change protocols.

\textbf{(3) Task-Coupled Consensus with Scheduler Replication.} We restrict consensus to dynamically formed task-coupled domains $C(\tau_j)$ of size $k \ll |N|$, reducing message complexity from $O(|N|^2)$ to $O(k^2)$. Scheduler replication ($m$ replicas) provides fault tolerance at the coordination layer; the specific fault model (Byzantine vs.\ crash) and replica count are clarified in Section~\ref{sec:protocol}.

\textbf{(4) Experimental Validation.} We validate STLC on a Raspberry Pi hardware testbed (6 agent nodes) and CORE network emulator (10--100 nodes). Hardware experiments confirm: 100\% detection of timestamp forgery and kinematic violations, 98\% detection of position spoofing (one false negative due to GPS noise), and Byzantine node isolation within 4--7 rounds. Emulation confirms 90\% early message filtering and 168\,ms consensus latency under 30\% Byzantine load, well within the 500\,ms real-time threshold.
%=============================================================================
\section{Related Work}
\label{sec:related}

Byzantine consensus research spans classical BFT protocols, modern optimizations, subset-based participation mechanisms, and application-specific designs. We focus here on identifying the gaps that motivate STLC's design; a quantitative performance comparison with representative baselines is deferred to Section~\ref{sec:experiments}.

\subsection{Classical BFT and Lightweight Mechanisms}

PBFT~\cite{habibi2023performance, ye2024tp} achieves strong consistency through three-phase message exchanges but requires all $|N|$ nodes to participate in every consensus round. Subsequent optimizations inherit this global participation requirement: SBFT~\cite{arun2022scalable} uses threshold signatures to reduce communication rounds; CheapBFT~\cite{wang2022bft} employs optimistic execution but reverts to PBFT-like complexity under faults; HotStuff~\cite{wang2024improved} achieves $O(n)$ per-round message complexity through leader-based pipelining but still operates globally across all validators.

Critically, \textit{none of these protocols validate the physical plausibility of messages}. They authenticate message origin via cryptographic signatures but do not check whether a reported position is kinematically reachable, whether a timestamp respects clock skew bounds, or whether a sender is legitimately assigned to the task in question. In industrial CPS, these physical invariants provide strong, $O(1)$-cost detection signals that can filter the majority of Byzantine messages before they consume consensus resources.

Recognizing IoT resource constraints, researchers developed lightweight adaptations. Wu et al.~\cite{wu2023learning} use reinforcement learning for adaptive protocol selection, while Nanapu et al.~\cite{nanapu2025d2bft} optimize PBFT parameters via deep RL. However, these methods address protocol \textit{tuning} rather than \textit{redesign}---the underlying protocols still assume global consensus and lack spatiotemporal validation mechanisms suited to cyber-physical environments.

\subsection{Byzantine Fault Tolerance in Industrial Cyber-Physical Systems}

Industrial multi-agent systems introduce unique challenges that classical BFT protocols do not address. Mii et al.~\cite{mivsic2020adapting} apply PBFT to industrial IoT networks but inherit its global participation requirement, making it impractical for large-scale factories. Leng et al.~\cite{leng2019manuchain} combine blockchain with multi-agent coordination for smart manufacturing, yet their consensus operates globally across all agents regardless of task locality. Recent work on resilient multi-agent systems~\cite{pirani2022survey} focuses on graph-theoretic conditions for consensus under Byzantine agents but assumes agents exchange states continuously without exploiting task-specific coupling or physical constraints.

Critically, none of these approaches leverage the \textit{physical measurability} of industrial agents---the fact that positions are kinematically bounded, velocities are motor-limited, and timestamps follow causal ordering---as a first-class mechanism for early Byzantine detection before entering the consensus protocol. This gap is the primary motivation for STLC's spatiotemporal validation layer.

\subsection{Trust-Based and Modern BFT Approaches}

Trust-based methods weight node influence by behavioral history. Karimireddy et al.~\cite{karimireddy2021learning} introduced learning-from-history for Byzantine-robust federated learning, while Mirzadeh et al.~\cite{mirzadeh2022filtering} proposed distributed malicious detection for vehicular networks. However, their trust models are \textit{static within consensus rounds}---scores update between instances but do not dynamically influence the ongoing consensus process. This creates a vulnerability window where Byzantine nodes can behave correctly to build trust, then inject malicious data before their scores are updated---problematic for industrial CPS where millisecond-scale tasks (e.g., AGV handoffs) require immediate Byzantine isolation to prevent physical collisions. STLC's adaptive trust weights, by contrast, are updated \textit{within} the consensus round based on spatiotemporal violation signals, enabling isolation within 4--7 rounds.

Recent DAG-based protocols achieve sublinear complexity through parallelism: Narwhal/Tusk~\cite{danezis2022narwhal} decouples data dissemination from consensus, achieving $O(1)$ amortized latency; Bullshark~\cite{spiegelman2022bullshark} improves with 2-round partially-synchronous consensus; Jolteon/Ditto~\cite{gelashvili2022jolteon} combines optimistic responsiveness with efficient view-change. These protocols optimize \textit{global consensus} through DAG parallelism but fundamentally assume all nodes must participate and do not exploit physical constraints for pre-consensus filtering. While effective for blockchain and permissionless systems, this global participation requirement remains impractical for industrial CPS where task-specific agent subsets (e.g., 3--5 agents per handoff) dominate operational patterns.

\subsection{Subset-Based and Weighted BFT Mechanisms}

A body of work reduces consensus participation to selected subsets of nodes, or introduces weighted voting to differentiate node influence. We survey three representative paradigms and identify their limitations for industrial CPS.

\textbf{Committee-Based BFT (Algorand~\cite{gilad2017algorand} and variants).} Algorand uses Verifiable Random Functions (VRFs) to cryptographically select a random committee of validators for each consensus round, achieving $O(1)$ per-round message complexity independent of total network size. This randomized selection provides strong unpredictability guarantees suitable for open, permissionless blockchains. However, random committee selection is fundamentally misaligned with industrial CPS requirements: task execution demands \textit{deterministic} participation by the specific agents physically involved in the task (e.g., the AGV transporting a component and the robot arm receiving it). A randomly selected committee of uninvolved agents cannot verify task-specific spatiotemporal claims---they have no access to the local sensor readings needed to check kinematic feasibility---nor can they execute the agreed decision. STLC instead forms consensus domains deterministically based on physical task coupling, providing both communication efficiency and semantic correctness of the decided value.

\textbf{Sharded BFT Systems (OmniLedger~\cite{kokoris2018omniledger}, RapidChain~\cite{zamani2018rapidchain}).} Sharding partitions the node set into fixed or epoch-based shards, each running independent consensus on a data partition. This achieves horizontal scalability by parallelizing consensus across shards. However, shard boundaries are determined by data partitioning criteria and remain stable across epochs, making them unsuitable for industrial environments where task-coupled agent sets change continuously as production workflows evolve. Furthermore, cross-shard transactions require expensive atomic commit protocols, whereas STLC's disjoint task-coupled domains are naturally independent and require no cross-domain coordination. Importantly, neither sharding nor committee-based approaches incorporate physical constraint validation: they reduce \textit{who} participates but do not filter \textit{what} messages are physically plausible before consensus.

\textbf{Weighted Byzantine Quorum Systems (Malkhi \& Reiter~\cite{malkhi1998byzantine}).} Weighted Byzantine quorum systems generalize classical quorum systems by assigning static weights to nodes, allowing quorum thresholds to reflect differential trust levels. However, weights are assigned at system initialization and remain fixed throughout operation. In industrial CPS, node trustworthiness changes dynamically: a previously honest AGV may become compromised mid-operation, or a transiently faulty sensor may recover. Static weights cannot capture this temporal evolution. STLC's adaptive trust management continuously updates weights based on real-time spatiotemporal violation signals, achieving Byzantine node isolation within 7 rounds without requiring explicit membership changes or system reconfiguration.

\subsection{Research Gaps}

\textbf{Gap 1: Absence of physical constraint exploitation for Byzantine detection.} All reviewed methods treat messages as abstract data and rely solely on cryptographic authentication. No prior work formalizes kinematic feasibility, temporal causality, and task-membership consistency as a pre-consensus Byzantine detection layer for industrial CPS. This represents a missed opportunity: physical laws provide $O(1)$-cost validation signals that can filter 65--90\% of Byzantine messages before they consume consensus resources, without cryptographic operations.

\textbf{Gap 2: Static or inter-round trust updates insufficient for millisecond-scale CPS.} Existing trust mechanisms update scores between consensus instances, creating a vulnerability window exploitable by Byzantine nodes that behave correctly to accumulate trust before injecting malicious data. No prior work integrates intra-round trust weight updates driven by spatiotemporal violation signals.

\textbf{Gap 3: Consensus globality and communication overhead.} All reviewed methods require global or quasi-global participation, resulting in $O(|N|^2)$ or $O(|N|)$ communication complexity even when tasks involve small agent subsets. While task-coupled consensus domains reduce communication, they do not by themselves guarantee equivalent global safety properties---a distinction we clarify in Section~\ref{sec:protocol} by scoping safety and liveness claims to the per-task domain level and providing a three-tier fallback mechanism for cases where the local Byzantine bound is violated.

\textbf{Gap 4: Static domains unable to adapt to dynamic task coupling.} Methods that reduce participation use static or slowly-changing agent sets determined at system initialization or epoch boundaries. Industrial tasks exhibit dynamic coupling where agents collaborate with different partners over time, requiring different consensus domains per task. Static approaches either include all potential collaborators (reverting to global consensus), or use fixed domains that cannot adapt when Byzantine attacks require expanding the validator set to maintain the $3f+1$ safety threshold.

\textbf{Gap 5: Scheduler-only consensus as an insufficient alternative.} A natural simplification is to run consensus exclusively on a small set of scheduler nodes, with all agents acting as passive clients---reducing participation to $O(m)$ nodes where $m \ll |N|$. While this reduces communication load, it introduces three critical limitations. \textit{First}, it creates a \textbf{trust centralization bottleneck}: all task execution decisions are determined solely by scheduler nodes, which have no direct physical involvement in tasks and cannot independently verify agent-reported spatiotemporal states (e.g., whether AGV$_1$ has actually arrived at the handoff position). \textit{Second}, \textbf{spatiotemporal validation is absent}: scheduler nodes are not co-located with task execution and cannot apply kinematic or spatial feasibility checks to agent messages, leaving the system vulnerable to position spoofing and timestamp forgery. \textit{Third}, \textbf{concurrent task parallelism is bottlenecked}: all tasks share the same scheduler consensus instance, creating a sequential throughput ceiling as task concurrency grows---whereas STLC's disjoint task-coupled domains execute fully in parallel with no inter-domain coordination. We include Scheduler-Only PBFT (S-PBFT) as an explicit baseline in our experimental evaluation (Section~\ref{sec:experiments}) to quantify these limitations empirically.

\section{System Model and Problem Formulation}
\label{sec:model}

We formalize the industrial multi-agent system, task-coupling relationships, Byzantine threat model, and optimization objectives.

\subsection{Agent and Communication Model}
\label{sec:agent_com_model}

\begin{definition}
The system consists of node set $N$ partitioned into three disjoint subsets: (i) \textit{Scheduler Set} $S$ containing $m$ coordinator replicas; (ii) \textit{AGV Set} containing $n_a$ mobile agents; (iii) \textit{Robot Arm Set} containing $n_r$ stationary agents. The total number of nodes is $|N| = m + n_a + n_r$.
\end{definition}

\textbf{Network Model.}
We adopt the \textit{partial synchrony} model~\cite{dwork1988consensus}. Specifically, the network operates in two phases:
\begin{itemize}
    \item \textit{Before GST (Global Stabilization Time)}: message delays are finite but unbounded. No liveness guarantees are made in this phase; safety (Agreement, Validity) holds regardless.
    \item \textit{After GST}: all message delays are bounded by a known constant $\Delta_{\max}$. Formally, any message sent at time $t \geq \text{GST}$ is guaranteed to be delivered by time $t + \Delta_{\max}$.
\end{itemize}

\textbf{Clarification on model terminology.}
We use ``partial synchrony'' throughout this paper to refer consistently to the Dwork--Lynch--Stockmeyer model above. The term ``asynchronous'' appearing in earlier drafts referred informally to the pre-GST phase only and has been removed to avoid ambiguity. Safety properties (Agreement, Validity, P1--P3) hold in \emph{both} phases. Liveness (Termination, P2) is guaranteed only after GST, within bounded time $T_{\max} = 9\Delta_{\max}$ (Theorem~\ref{thm:liveness}).

\textbf{Communication Model.}
Agents communicate over network graph $\mathcal{G} = (N, \mathcal{E})$, where $\mathcal{E}$ denotes the set of communication links. Each message $\mathbf{m}$ has structure:
\begin{equation}
\begin{aligned}
\mathbf{m} = \langle \text{id}_{\text{sender}},\; t_{\text{stamp}},\; \mathbf{p}_{\text{loc}},\; \text{id}_{\text{task}},\; \text{payload} \rangle
\end{aligned}
\end{equation}
where $\text{id}_{\text{sender}} \in N$ identifies the sender, $t_{\text{stamp}} \in \mathbb{R}_{\geq 0}$ is the timestamp, $\mathbf{p}_{\text{loc}} \in \mathbb{R}^2$ is the claimed position, $\text{id}_{\text{task}}$ identifies the associated task, and $\text{payload}$ contains task-specific data. All messages are authenticated via digital signatures; the adversary cannot forge signatures of honest nodes (Assumption~\ref{assump:crypto}).

\textbf{Spatiotemporal Model.}
The system operates in continuous time domain $\mathbb{R}_{\geq 0}$ with a global reference clock. Each node $n_i \in N$ maintains a local clock with bounded drift $\delta_t$ from the global clock:
\[
    |t_{\text{local}}^{(i)} - t_{\text{global}}| \leq \delta_t,
    \quad \forall\, t_{\text{global}} \in \mathbb{R}_{\geq 0}
\]
Agents operate in a 2D Euclidean workspace $\mathbb{R}^2$. The position of node $n_i$ at time $t$ is denoted $\mathbf{p}(n_i, t) \in \mathbb{R}^2$. AGVs are mobile with maximum velocity $v_{\max} > 0$, satisfying the kinematic constraint:
\begin{equation}
\begin{aligned}
\|\mathbf{p}(n_i, t_2) - \mathbf{p}(n_i, t_1)\| \leq
v_{\max} \cdot (t_2 - t_1), \\
\quad \forall\, n_i \in \text{AGV Set},\;
\forall\, t_1, t_2 \in \mathbb{R}_{\geq 0}
\text{ with } t_2 \geq t_1
\end{aligned}
\end{equation}
where $\|\cdot\|$ denotes the Euclidean norm. Robot Arms are stationary: $\mathbf{p}(n_i, t) = \mathbf{p}(n_i, 0)$ for all $t$.

\subsection{Task Coupling Model}
\label{sec:task_model}

\begin{definition}[Task]
\label{def:task}
A task $\tau_j$ is a tuple $\langle \text{id}_\tau,\; \text{type}_\tau,\; \mathcal{A}_{\text{req}},\; \mathcal{T}_{\text{win}},\; \mathcal{P}_{\text{reg}} \rangle$ where:
\begin{itemize}
    \item $\text{id}_\tau \in \mathbb{N}$ is the unique task identifier,
    \item $\text{type}_\tau$ specifies the task category (e.g., material handoff, assembly),
    \item $\mathcal{A}_{\text{req}} \subseteq N$ is the set of required agent types,
    \item $\mathcal{T}_{\text{win}} = [t_{\text{start}}, t_{\text{end}}] \subset \mathbb{R}_{\geq 0}$ is the temporal execution window with $t_{\text{start}} < t_{\text{end}}$,
    \item $\mathcal{P}_{\text{reg}} \subset \mathbb{R}^2$ is the spatial region where the task is executed.
\end{itemize}
\end{definition}

\begin{definition}[Task-Coupled Node Set]
\label{def:task_coupled_set}
For task $\tau_j$, the \textit{task-coupled set} $C(\tau_j) \subseteq N$ consists of all nodes involved in executing or coordinating $\tau_j$:
\begin{equation}
\begin{aligned}
C(\tau_j) = S \cup \{n_i \in N \mid n_i \text{ is assigned to execute } \tau_j\}
\end{aligned}
\end{equation}
where $S$ is the scheduler set, and the second term includes all agents directly participating in task execution. We denote $k_j = |C(\tau_j)|$ as the size of the task-coupled set for task $\tau_j$.
\end{definition}

\textbf{Communication Complexity Reduction.}
By restricting consensus to the task-coupled set $C(\tau_j)$ rather than the entire node set $N$, message complexity reduces from $O(|N|^2)$ to $O(k_j^2 + m^2 + mk_j)$, where $k_j \ll |N|$ and $m = |S|$. This reduction is a secondary benefit of task-coupled consensus; the primary motivation is that only task-coupled agents possess the local sensor readings required to apply spatiotemporal validation (Section~\ref{sec:protocol}). Agents outside $C(\tau_j)$ cannot verify kinematic feasibility of task-specific position reports and therefore cannot meaningfully participate in the validation layer.

\subsection{Byzantine Threat Model}
\label{sec:threat_model}

\begin{definition}[Adversary Model]
\label{def:adversary}
Let $\mathcal{F} \subseteq N$ denote the set of Byzantine (compromised) nodes with $|\mathcal{F}| = f$. The adversary has the following capabilities:
\begin{enumerate}
    \item \textit{Message forgery}: Create messages with false timestamps $t_{\text{stamp}}$ or false positions $\mathbf{p}_{\text{loc}}$,
    \item \textit{Message replay}: Resend previously valid messages at different times,
    \item \textit{Selective participation}: Choose whether to respond to consensus requests,
    \item \textit{Collusion}: Coordinate malicious behavior among compromised nodes in $\mathcal{F}$.
\end{enumerate}
The adversary \textit{cannot}: (a) break cryptographic primitives (e.g., digital signatures), or (b) compromise more than $f < |N|/3$ nodes globally.
\end{definition}

\begin{assumption}[Cryptographic Integrity]
\label{assump:crypto}
All honest nodes sign messages with unforgeable digital signatures (e.g., Ed25519). A message bearing a valid signature of node $n_i$ was sent by $n_i$.
\end{assumption}

\begin{assumption}[Failure Detection: Completeness]
\label{assump:detection_completeness}
Spatiotemporal detection is \emph{not} complete in the strict sense: a Byzantine node $n_i \in \mathcal{F}$ that crafts messages within sensor noise bounds (i.e., position error $\leq \varepsilon_s$ and timestamp error $\leq \delta_t$) may evade detection in a given round. We distinguish two cases:

\begin{itemize}
    \item \textit{Persistent violators}: A Byzantine node that sends at least one detectable violation per consensus round is detected within $K$ rounds, where $K$ is the weight convergence upper bound implied by Eq.~(\ref{eq:isolation_time}) (experimentally $K \leq 10$ rounds). Its adaptive weight is penalized by factor $\rho_p < 1$ upon each detected violation, reducing quorum influence even before full isolation.
    \item \textit{Stealthy violators}: A Byzantine node that consistently stays within sensor noise bounds can evade spatiotemporal detection entirely. Such nodes are handled by the BFT quorum layer: even if a stealthy Byzantine node participates in consensus with full weight, the quorum requirement $f_j < k_j/3$ ensures it cannot unilaterally determine the decided value. Safety is preserved by the quorum intersection property (Theorem~\ref{thm:safety}), not by detection completeness.
\end{itemize}

The false negative rate for single-round violations is bounded by $P_{\text{fn}} = P(\|\mathbf{e}_s\| > \varepsilon_s) + P(|e_t| > \delta_t)$ under the sensor error bounds of Assumption~\ref{assump:sensor_accuracy}.
\end{assumption}

\begin{assumption}[Failure Detection: Accuracy]
\label{assump:sensor_accuracy}
Correct nodes' sensors have bounded measurement errors: position error $\|\mathbf{e}_s\| \leq \varepsilon_s = 0.5\,\text{m}$ and clock drift $|e_t| \leq \delta_t = 100\,\text{ms}$. The spatiotemporal validation thresholds are set to strictly subsume these error bounds:
\[
  \epsilon_{\text{spatial}} = \varepsilon_s = 0.5\,\text{m},
  \qquad
  \Delta_t = \delta_t + \Delta_{\max}
\]
Consequently, a message from an honest node with bounded sensor noise always passes validation---the system produces \emph{no false positives} against honest nodes under this assumption.

\textbf{Safety under sensor noise.} Safety is preserved even when sensor noise is present because consensus decisions are derived from cross-validated reports of an honest quorum, not from any single agent's self-report. Specifically, the decided value $v^*$ requires agreement from at least $\lfloor 2k_j/3 \rfloor + 1$ agents; a single honest agent's noisy report can shift $v^*$ by at most $\varepsilon_s$ from the true value, which is within the operational tolerance of industrial handoff tasks (position accuracy requirement: $\pm 1\,\text{m}$).
\end{assumption}

\begin{assumption}[Localized Byzantine Bound]
\label{assump:localized_bound}
For task $\tau_j$ with task-coupled set $C(\tau_j)$ of size $k_j = |C(\tau_j)|$, let $f_j = |\mathcal{F} \cap C(\tau_j)|$ denote the number of Byzantine nodes within the task-coupled domain. We require:
\begin{equation}
\begin{aligned}
f_j < \frac{k_j}{3}
\end{aligned}
\end{equation}
If this bound is violated in the initial domain $C(\tau_j)$, adaptive domain expansion (Section~\ref{sec:adaptive_expansion}) is triggered to restore the safety margin.
\end{assumption}

\begin{assumption}[Byzantine Distribution]
\label{assump:byz_distribution}
The theoretical safety and complexity analyses in Theorems~\ref{thm:safety}--\ref{thm:complexity} assume a \emph{worst-case uniform random} distribution of Byzantine nodes across $N$: each node is independently compromised with probability $f/|N|$. This is the standard assumption in BFT literature and yields the hypergeometric violation probability analysis of Equation~\eqref{eq:violation_prob}.

We explicitly acknowledge that in industrial CPS, Byzantine failures may be \emph{spatially clustered}: compromised devices often share common software vulnerabilities or are targeted by an attacker focused on a high-value task region. Such clustering \emph{increases} the probability that a single task-coupled domain $C(\tau_j)$ contains multiple Byzantine nodes, potentially violating Assumption~\ref{assump:localized_bound}. STLC addresses clustered Byzantine attacks through: (1)~\textit{Spatiotemporal validation} (Layer~1) detecting and rejecting physically inconsistent messages regardless of clustering; (2)~\textit{Adaptive domain expansion} (Section~\ref{sec:adaptive_expansion}) restoring the $f_j < k_j/3$ threshold by including additional honest nodes. Experimental evaluation of clustered attacks (Section~\ref{sec:clustered}) confirms 100\% task completion under clustered scenarios with $f_j/k_j \leq 1/3 + \epsilon$ via expansion, with 98.7\% overall success rate.
\end{assumption}

\begin{assumption}[Weight Update Consistency]
\label{assump:weight_consistency}
Adaptive weight updates and domain expansion decisions are \emph{not} required to be globally consistent across all replicas in real time. Each agent in $C(\tau_j)$ independently computes weight updates based on locally observed spatiotemporal violations. Safety is preserved because the quorum threshold is evaluated against the \emph{local} weight vector of each agent: two honest agents with slightly different weight vectors may have different quorum thresholds, but the quorum intersection property (Theorem~\ref{thm:safety}) holds as long as both weight vectors satisfy $W_{\text{Byzantine}} < W_{\text{total}}/3$, which is guaranteed by the weight update bounds in Section~\ref{sec:layer3}. Liveness is preserved because weight convergence occurs within $K$ rounds, after which all honest agents' weight vectors agree on the isolation of persistent Byzantine nodes.
\end{assumption}

\textbf{Attack Scenarios.}
We defend against three primary attack vectors:
\begin{enumerate}
    \item \textit{Timestamp forgery}: Messages with timestamps $t_{\text{stamp}} \notin [t_{\text{current}} - \Delta_t,\; t_{\text{current}} + \Delta_t]$ indicate replay attacks or future-dated messages, where $\Delta_t = \delta_t + \Delta_{\max}$,
    \item \textit{Spatial forgery}: Reported positions $\mathbf{p}_{\text{loc}}$ violating kinematic constraints, i.e., $\|\mathbf{p}_{\text{loc}} - \mathbf{p}_{\text{expected}}\| > \epsilon_{\text{spatial}}$,
    \item \textit{Task inconsistency}: Nodes claiming involvement in task $\tau_j$ when they are not members of $C(\tau_j)$.
\end{enumerate}

\textbf{Probabilistic Analysis of Assumption~\ref{assump:localized_bound}.}
For a task-coupled domain $C(\tau_j)$ of size $k_j$ selected from $|N|$ total nodes containing $f$ Byzantine nodes under uniform random distribution, the probability that $f_j \geq \lceil k_j/3 \rceil$ follows a hypergeometric distribution:
\begin{equation}
\begin{aligned}
P\!\left(f_j \geq \left\lceil \frac{k_j}{3} \right\rceil\right) =
\sum_{i=\lceil k_j/3 \rceil}^{\min(f,\,k_j)}
\frac{\dbinom{f}{i}\dbinom{|N|-f}{k_j-i}}{\dbinom{|N|}{k_j}}
\label{eq:violation_prob}
\end{aligned}
\end{equation}

Table~\ref{tab:violation_prob} shows violation probabilities for typical industrial configurations.

\begin{table}[h]
\centering
\caption{Violation Probability of Assumption~\ref{assump:localized_bound} under Uniform Random Byzantine Distribution}
\label{tab:violation_prob}
\small
\begin{tabular}{ccccc}
\toprule
$k_j$ & $|N|$ & $f$ & $P(f_j \geq k_j/3)$ & \textbf{Risk Level}$^{\dagger}$ \\
\midrule
3  & 30 & 10 & 0.716 (71.6\%) & \textcolor{red}{HIGH} \\
3  & 30 & 5  & 0.386 (38.6\%) & \textcolor{orange}{MEDIUM} \\
7  & 30 & 10 & 0.182 (18.2\%) & \textcolor{orange}{MEDIUM} \\
7  & 30 & 5  & 0.031 (3.1\%)  & \textcolor{green}{LOW} \\
10 & 30 & 10 & 0.051 (5.1\%)  & \textcolor{green}{LOW} \\
\bottomrule
\end{tabular}
\vspace{0.2cm}

\noindent $^{\dagger}$Risk levels defined according to IEC 61508 Safety Integrity Levels~\cite{iec61508}: HIGH ($P > 10\%$, exceeds SIL~1 upper bound), MEDIUM ($1\% < P \leq 10\%$, SIL~1 range), LOW ($P \leq 1\%$, SIL~2 or higher).
\end{table}

\textbf{Practical Justification for $k_j = 3$.}
Table~\ref{tab:violation_prob} shows $k_j = 3$ has 71.6\% violation probability under worst-case uniform random Byzantine distribution ($f/|N| = 1/3$). Real-world deployments exhibit three mitigating factors: (1)~\textit{Spatial clustering} of Byzantine nodes reduces violation probability to 12.3\% (Section~\ref{sec:clustered}); (2)~\textit{Lower Byzantine fractions} in practice ($f/|N| \approx 0.16$ vs.\ worst-case 0.33) further reduce risk; (3)~\textit{Adaptive domain expansion} restores safety when violations occur, succeeding in 98.7\% of cases (Section~\ref{sec:expansion_eval}). For safety-critical applications, we recommend $k_j \geq 7$ (violation probability $< 20\%$) to balance efficiency and fault tolerance.

\subsection{Problem Formulation}
\label{sec:problem_formulation}

\begin{problem}[Efficient Byzantine-Resilient Consensus for Task-Coupled Multi-Agent Systems]
\label{prob:main}
\textbf{Given:}
\begin{itemize}
    \item A multi-agent system with node set $N$, communication graph $\mathcal{G} = (N, \mathcal{E})$, and time domain $\mathbb{R}_{\geq 0}$ (Section~\ref{sec:agent_com_model}),
    \item A set of tasks $\mathcal{T} = \{\tau_1, \tau_2, \ldots, \tau_{|\mathcal{T}|}\}$ with task-coupled sets $\{C(\tau_j)\}_{j=1}^{|\mathcal{T}|}$ (Definition~\ref{def:task_coupled_set}),
    \item A set of Byzantine nodes $\mathcal{F} \subseteq N$ with $|\mathcal{F}| = f < |N|/3$ (Definition~\ref{def:adversary}),
\end{itemize}

\noindent \textbf{Design} a consensus protocol that satisfies the following properties:

\vspace{0.3cm}
\noindent \textbf{P1. Safety (Agreement):} For any task $\tau_j \in \mathcal{T}$, all honest nodes in the task-coupled set $C(\tau_j)$ agree on the same decision value $v_j$:
\begin{equation}
\begin{aligned}
\forall\, n_i, n_k \in (C(\tau_j) \setminus \mathcal{F}):\quad
\text{decision}(n_i, \tau_j) = \text{decision}(n_k, \tau_j) = v_j
\label{eq:safety}
\end{aligned}
\end{equation}
where $\text{decision}(n_i, \tau_j)$ denotes the final decision of node $n_i$ for task $\tau_j$.

\textbf{Scope of safety.} Safety is defined \emph{per task domain} $C(\tau_j)$, not globally across all $N$ nodes. Nodes outside $C(\tau_j)$ do not participate in and are not required to agree on $v_j$. System-wide consistency is achieved through the scheduler layer, which assigns globally ordered sequence numbers to tasks before domain broadcast (Remark~\ref{rem:ordering}).

\vspace{0.3cm}
\noindent \textbf{P2. Liveness (Termination):} For any task $\tau_j \in \mathcal{T}$ initiated at time $t_{\text{init}}$ \emph{after GST}, all honest nodes in $C(\tau_j)$ eventually reach a decision within bounded time $T_{\max}$:
\begin{equation}
\begin{aligned}
\exists\, t \in [t_{\text{init}},\; t_{\text{init}} + T_{\max}]:\;
\forall\, n_i \in (C(\tau_j) \setminus \mathcal{F}),\;
\text{decision}(n_i, \tau_j) \neq \bot
\label{eq:liveness}
\end{aligned}
\end{equation}
where $\bot$ denotes the undecided state, and $T_{\max} = 9\Delta_{\max}$ under partial synchrony after GST (Theorem~\ref{thm:liveness}). Before GST, liveness is not guaranteed; safety holds unconditionally.

\vspace{0.3cm}
\noindent \textbf{P3. Validity:} The agreed decision value $v_j$ must be proposed by at least one honest node in $C(\tau_j)$:
\begin{equation}
\begin{aligned}
\exists\, n_i \in (C(\tau_j) \setminus \mathcal{F}):\;
\text{proposal}(n_i, \tau_j) = v_j
\label{eq:validity}
\end{aligned}
\end{equation}
where $\text{proposal}(n_i, \tau_j)$ is the initial proposal of node $n_i$ for task $\tau_j$. This prevents Byzantine nodes from forcing arbitrary decisions.

\vspace{0.3cm}
\noindent \textbf{P4. Byzantine Tolerance:} Properties P1--P3 hold under the following conditions:
\begin{itemize}
    \item \textit{Global bound}: $f < |N|/3$ (guaranteed by Definition~\ref{def:adversary}),
    \item \textit{Local bound}: For each task $\tau_j$, the protocol ensures $f_j < |C'(\tau_j)|/3$ where $C'(\tau_j)$ is the (possibly expanded) consensus domain. If the initial domain $C(\tau_j)$ violates this bound, adaptive domain expansion is triggered to restore the safety margin.
\end{itemize}

\vspace{0.3cm}
\noindent \textbf{P5. Task Locality:} For each task $\tau_j$, consensus involves only nodes in the task-coupled set $C(\tau_j)$ (or its expanded version $C'(\tau_j)$), not the entire node set $N$:
\begin{equation}
\begin{aligned}
\forall\, n_i \in (N \setminus C(\tau_j)):\;
\text{participates}(n_i, \tau_j) = \textsc{False}
\label{eq:locality}
\end{aligned}
\end{equation}
where $\text{participates}(n_i, \tau_j)$ indicates whether node $n_i$ sends or receives consensus messages for task $\tau_j$.

\vspace{0.3cm}
\noindent \textbf{Optimization Objective:} Among all protocols satisfying P1--P5, minimize the weighted sum of communication cost and latency:
\begin{equation}
\begin{aligned}
\min \quad \alpha \cdot \mathbb{E}[\text{CommCost}] +
\beta \cdot \mathbb{E}[\text{Latency}]
\label{eq:objective}
\end{aligned}
\end{equation}
subject to: Safety, Liveness, Validity, Byzantine Tolerance, Task Locality (P1--P5), where:
\begin{itemize}
    \item $\text{CommCost} = \sum_{\tau_j \in \mathcal{T}} M_j$ is the total number of messages, with $M_j$ denoting messages sent during consensus for task $\tau_j$,
    \item $\text{Latency} = \max_{\tau_j \in \mathcal{T}} (t_{\text{decide}}^{(j)} - t_{\text{init}}^{(j)})$ is the maximum time from task initiation to decision,
    \item $\mathbb{E}[\cdot]$ denotes expectation over task distributions and Byzantine behaviors,
    \item $\alpha, \beta \geq 0$ are application-specific weights with $\alpha + \beta = 1$.
\end{itemize}
\end{problem}

\vspace{0.3cm}
\noindent \textbf{Remark 1 (Single-Shot Consensus per Task and Decision Order).}
\label{rem:ordering}
Each task $\tau_j$ requires exactly \emph{one} consensus decision---a single-shot consensus instance producing a decided value:
\[
  v_j^* = \langle p^*_{\text{AGV}},\; t^*_{\text{handoff}},\; \text{id}^*_{\text{component}} \rangle
\]
representing the cross-validated, agreed task execution parameters. Multi-phase tasks decompose into independent single-shot consensus instances, one per phase, each running a fresh STLC instance.

Regarding decision ordering: decisions for tasks in \emph{disjoint} domains ($C(\tau_j) \cap C(\tau_{j'}) = \emptyset$) are naturally independent and require no global ordering. For tasks sharing agents (overlapping domains), the scheduler layer assigns monotonically increasing global sequence numbers $\text{seq}(\tau_j)$ before domain broadcast, ensuring agents participating in multiple tasks process decisions in a consistent order. This sequence number is included in the PRE-PREPARE message and verified by all domain members before entering the prepare phase.

\vspace{0.3cm}
\noindent \textbf{Remark 2 (Adaptive Byzantine Tolerance).}
Unlike static PBFT protocols that assume a fixed participant set, our formulation allows \textit{adaptive domain expansion}: if the initial task-coupled set $C(\tau_j)$ violates the local Byzantine bound $f_j < |C(\tau_j)|/3$, the protocol dynamically expands to $C'(\tau_j)$ to restore the safety margin (Section~\ref{sec:adaptive_expansion}). Experimental results (Section~\ref{sec:clustered}) show that expansion is triggered in only 8.3\% of tasks under realistic attack scenarios, with 98.7\% success rate.

%=============================================================================
\section{STLC: Spatiotemporal-First Byzantine Consensus for Industrial Multi-Agent Systems}
\label{sec:protocol}
\label{sec:protocol_design}

STLC is a Byzantine-resilient consensus protocol designed for industrial multi-agent systems with replicated scheduling infrastructure. Its central design principle is \textit{spatiotemporal-first}: physical constraints (kinematic bounds, clock causality, task membership) are enforced as a pre-consensus filter, rejecting 65--90\% of Byzantine messages at $O(1)$ cost before they consume consensus resources. The remaining consensus protocol is then restricted to task-coupled domains of size $k \ll |N|$, achieving $O(k^2)$ communication complexity. Safety and liveness are defined \emph{per task domain} $C(\tau_j)$ under the localized Byzantine bound $f_j < k_j/3$ (Assumption~\ref{assump:localized_bound}); a three-tier fallback mechanism handles cases where this bound is violated.

\subsection{Protocol Overview and Three-Layer Architecture}
\label{sec:protocol_overview}

STLC comprises three layers that work synergistically (Figure~\ref{fig:architecture}).

\begin{figure*}[!t]
\centering
\includegraphics[width=.85\textwidth]{fig/architecture_and_flow.pdf}
\Description{Three-layer STLC architecture showing Layer 1 spatiotemporal validation, Layer 2 task-coupled consensus with weighted quorums and view-change, and Layer 3 adaptive trust management.}
\caption{STLC three-layer architecture with message flow for task $\tau_1$ with $k=3$. Layer~1 performs $O(1)$ spatiotemporal validation to filter Byzantine messages before consensus. Layer~2 executes task-coupled consensus within localized domain $C(\tau_1)$ using weighted quorums and an integrated view-change sub-protocol. Layer~3 maintains adaptive trust weights that evolve based on node behavior, serving as a second line of defense against stealthy Byzantine nodes.}
\label{fig:architecture}
\end{figure*}

\textbf{Layer 1: Spatiotemporal Validation} (Section~\ref{sec:layer1}) performs $O(1)$ early rejection of physically inconsistent messages using embedded timestamps $t_{\text{msg}}$ and spatial coordinates $\mathbf{p}_{\text{msg}}$. Messages violating kinematic constraints (Equation~\ref{eq:kinematic_validation}) or temporal bounds (Equation~\ref{eq:temporal_validation}) are immediately discarded \textit{before} entering the consensus protocol. This layer exploits the physical measurability of industrial agents---a property absent in abstract distributed systems---to provide strong Byzantine detection at negligible cost.

\textbf{Layer 2: Task-Coupled Consensus} (Section~\ref{sec:layer2}) restricts consensus to dynamically formed domains $C(\tau_j)$ containing only task-relevant nodes. Consensus executes through three phases (pre-prepare, prepare, commit) \textit{with an integrated view-change sub-protocol} (Algorithm~\ref{alg:view_change}) among $k$ nodes, reducing message complexity from $O(N^2)$ to $O(k^2)$. The view-change sub-protocol is essential for liveness: without it, a Byzantine or crashed coordinator would halt consensus indefinitely. The layer also includes adaptive domain expansion and a three-tier fallback mechanism to handle cases where the localized bound $f_{\text{local}} < k/3$ is violated.

\textbf{Layer 3: Adaptive Trust Management} (Section~\ref{sec:layer3}) maintains dynamic weights $w_i$ for each node based on behavioral history. Nodes exhibiting suspicious activities (detected by Layer~1 or Layer~2) have weights reduced via multiplicative penalty. Consensus employs weighted voting, enabling gradual isolation of malicious nodes without requiring explicit Byzantine node identification. Weight updates are computed independently by each honest node based on locally observed violations; safety under bounded weight divergence between replicas is guaranteed by the quorum intersection argument (Theorem~\ref{thm:safety} and Assumption~\ref{assump:weight_consistency}).

This layered design ensures:
\begin{itemize}
\item \textbf{Early filtering} (Layer~1) reduces consensus load by rejecting physically impossible messages at $O(1)$ cost,
\item \textbf{True consensus} (Layer~2) with view-change guarantees both safety and liveness under partial synchrony,
\item \textbf{Adaptive isolation} (Layer~3) provides a second line of defense against stealthy Byzantine nodes that evade Layer~1 by staying within sensor noise bounds.
\end{itemize}

\subsubsection{Layer 1: Spatiotemporal Validation}
\label{sec:layer1}

Layer~1 performs $O(1)$ early validation across three dimensions before messages enter the consensus protocol, filtering 65--90\% of Byzantine messages (Table~\ref{tab:layer1_rejection}).

\paragraph{Validation Dimensions}

Each consensus message $\mathbf{m}$ is validated along three dimensions:

\textbf{(1) Temporal Validation.}
Messages with timestamp $t_{\text{msg}}$ are accepted only if:
\begin{equation}
\begin{aligned}
|t_{\text{recv}} - t_{\text{msg}}| \leq \Delta t_{\max}
\label{eq:temporal_validation}
\end{aligned}
\end{equation}
where $\Delta t_{\max} = \delta_t + \Delta_{\max} = 100\,\text{ms} + 400\,\text{ms} = 500\,\text{ms}$ accounts for clock drift (Assumption~\ref{assump:sensor_accuracy}) and network delay after GST. This detects replay attacks and timestamp forgery.

\textbf{(2) Spatial Validation.}
Messages with claimed position $\mathbf{p}_{\text{msg}}$ must satisfy kinematic constraints:
\begin{equation}
\begin{aligned}
\|\mathbf{p}_{\text{msg}} - \mathbf{p}_{\text{last}}\|
\leq v_{\max} \cdot (t_{\text{msg}} - t_{\text{last}}) +
\varepsilon_s
\label{eq:kinematic_validation}
\end{aligned}
\end{equation}
where $\mathbf{p}_{\text{last}}$ is the sender's previous position, $v_{\max} = 2\,\text{m/s}$ is the maximum velocity, and $\varepsilon_s = 0.5\,\text{m}$ is the sensor position error bound (Assumption~\ref{assump:sensor_accuracy}). This prevents position spoofing and teleportation attacks.

\textbf{(3) Task Consistency Validation.}
Messages are accepted only if the sender is authorized for task $\tau_j$:
\begin{equation}
\begin{aligned}
\text{sender}(\mathbf{m}) \in C(\tau_j)
\label{eq:task_consistency}
\end{aligned}
\end{equation}
This prevents cross-task injection attacks.

\paragraph{False Positive Analysis}

A \textit{false positive} occurs when a message from an honest node is incorrectly rejected by Layer~1. By Assumption~\ref{assump:sensor_accuracy}, honest nodes have bounded sensor errors: $\|\mathbf{e}_s\| \leq \varepsilon_s$ and $|e_t| \leq \delta_t$. The validation thresholds are designed with explicit safety margins:
\begin{itemize}
    \item \textit{Temporal threshold}: $\Delta t_{\max} = \delta_t + \Delta_{\max} = 500\,\text{ms}$ strictly subsumes the maximum honest clock drift plus post-GST network delay. An honest node's message always satisfies $|t_{\text{recv}} - t_{\text{msg}}| \leq \delta_t + \Delta_{\max} = \Delta t_{\max}$.
    \item \textit{Spatial threshold}: The additive margin $\varepsilon_s$ in Equation~\eqref{eq:kinematic_validation} subsumes GPS measurement noise. An honest node whose true position satisfies the kinematic constraint will always have its reported position within the threshold.
\end{itemize}
Therefore, \textbf{Layer~1 produces zero false positives} against honest nodes under Assumption~\ref{assump:sensor_accuracy}, confirmed experimentally: 0\% false positive rate across 150 controlled injection trials (Section~\ref{sec:testbed_validation}).

\paragraph{False Negative Analysis}

A \textit{false negative} occurs when a Byzantine message passes Layer~1. Two scenarios arise:

\textbf{(1) Subtle violations:} A Byzantine node reports a position kinematically feasible but strategically incorrect (e.g., $0.3\,\text{m}$ offset, within $\varepsilon_s$). The falsification magnitude is bounded by $\varepsilon_s$, so the resulting consensus decision deviates from ground truth by at most $\varepsilon_s$---within the physical tolerance of industrial handoff tasks.

\textbf{(2) Coordinated mimicry:} A Byzantine node temporarily behaves correctly to build trust before injecting malicious data. This is the primary false negative risk and is addressed by Layer~3 (see below).

\textbf{Layer~3 as second line of defense:} Both scenarios are addressed by Layer~3's adaptive weight mechanism. Even if a Byzantine message passes Layer~1 in a given round, the node's subsequent behavior (conflicting PREPARE messages, inconsistent commits) triggers weight penalties. A node that injects malicious data after a period of correct behavior will have its weight reduced upon detection, and its influence on future quorums diminishes within $T_{\text{isolate}} = 9$ rounds (Equation~\ref{eq:isolation_time}). Critically, \textit{no single false negative can compromise safety}: the weighted quorum threshold (Equation~\ref{eq:weighted_quorum}) requires $\sum_{n_k \in \mathcal{P}} w_k \geq \frac{2}{3} W_{\text{total}}$, so a single Byzantine message cannot satisfy the quorum without honest node participation.

\paragraph{Computational Complexity}

Each validation requires $O(1)$ operations (arithmetic comparison, vector norm, hash table lookup). Table~\ref{tab:layer1_rejection} shows early rejection rates under various attacks.

\begin{table}[h]
\centering
\caption{Layer 1 Early Rejection Rate Under Byzantine Attacks}
\label{tab:layer1_rejection}
\small
\adjustbox{width=\linewidth}{
\begin{tabular}{lcc}
\toprule
\textbf{Attack Type} & \textbf{CORE} & \textbf{Testbed} \\
\midrule
Timestamp forgery              & 95\% & 90\% \\
Position spoofing              & 88\% & 82\% \\
Task injection                 & 100\% & 100\% \\
Mixed attacks (30\% Byzantine) & 78\% & 65\% \\
\bottomrule
\end{tabular}}
\end{table}

\begin{algorithm}[t]
\caption{Weighted Three-Phase Consensus (Normal Path)}
\label{alg:weighted_consensus}
\small
\begin{algorithmic}[1]
\STATE \textbf{Input:} Task $\tau_j$, Domain $C(\tau_j)$,
       Weights $\{w_i\}$, View $v \gets 0$
\STATE \textbf{Output:} Decision $v_j$ for task $\tau_j$

\STATE \textbf{// Phase 1: Pre-Prepare (Coordinator $n_0$)}
\IF{$\text{isCoordinator}(v, C(\tau_j))$}
  \STATE Compute $d(\tau_j) = \text{Hash}(\tau_j)$
  \STATE Broadcast
         $\langle \text{PRE-PREPARE},\, v,\, s,\, d(\tau_j)
         \rangle_{\sigma_{n_0}}$ to $C(\tau_j)$
\ENDIF

\STATE \textbf{// Phase 2: Prepare}
\UPON{receive $\langle \text{PRE-PREPARE},\, v,\, s,\, d \rangle$
      from coordinator}
  \IF{$\text{ValidateSpatiotemporal}(\mathbf{m}) \land
      \text{VerifySignature}(\mathbf{m})$}
    \STATE Start timer $\mathcal{T}_{\text{vc}}$ with timeout
           $\tau_{\text{to}} = 3\Delta_{\max}$
    \STATE Broadcast
           $\langle \text{PREPARE},\, v,\, s,\, d
           \rangle_{\sigma_{\text{self}}}$ to $C(\tau_j)$
    \STATE $\mathcal{P} \gets \{\text{coordinator}\}$
  \ENDIF
\ENDUPON

\UPON{receive $\langle \text{PREPARE},\, v,\, s,\, d \rangle$
      from $n_k$}
  \IF{$\text{ValidateSpatiotemporal}(\mathbf{m}) \land
      \text{VerifySignature}(\mathbf{m})$}
    \STATE $\mathcal{P} \gets \mathcal{P} \cup \{n_k\}$
    \STATE $W_{\mathcal{P}} \gets \sum_{n_i \in \mathcal{P}} w_i$;
           \quad $W_{\text{total}} \gets
           \sum_{n_i \in C(\tau_j)} w_i$
    \IF{$W_{\mathcal{P}} \geq \frac{2}{3} W_{\text{total}}$}
      \STATE Cancel $\mathcal{T}_{\text{vc}}$
      \STATE Broadcast
             $\langle \text{COMMIT},\, v,\, s,\, d
             \rangle_{\sigma_{\text{self}}}$ to $C(\tau_j)$
      \STATE $\mathcal{C} \gets \{\text{self}\}$
    \ENDIF
  \ENDIF
\ENDUPON

\STATE \textbf{// Phase 3: Commit}
\UPON{receive $\langle \text{COMMIT},\, v,\, s,\, d \rangle$
      from $n_k$}
  \IF{$\text{ValidateSpatiotemporal}(\mathbf{m}) \land
      \text{VerifySignature}(\mathbf{m})$}
    \STATE $\mathcal{C} \gets \mathcal{C} \cup \{n_k\}$
    \STATE $W_{\mathcal{C}} \gets \sum_{n_i \in \mathcal{C}} w_i$
    \IF{$W_{\mathcal{C}} \geq \frac{2}{3} W_{\text{total}}$}
      \STATE $v_j \gets \text{ExecuteDecision}(\tau_j,\, d)$
      \STATE $\text{UpdateWeights}(C(\tau_j),\, \mathcal{P},\,
             \mathcal{C})$ \COMMENT{Layer~3}
      \RETURN $v_j$
    \ENDIF
  \ENDIF
\ENDUPON

\STATE \textbf{// View-Change (on timer expiry)}
\UPON{$\mathcal{T}_{\text{vc}}$ expires in view $v$}
  \STATE \textbf{goto} Algorithm~\ref{alg:view_change}, line~1
\ENDUPON
\end{algorithmic}
\end{algorithm}

\paragraph{View-Change Sub-Protocol}
\label{sec:view_change}

The view-change sub-protocol ensures liveness when the current coordinator fails or behaves Byzantine, and is what distinguishes STLC from a reliable broadcast protocol.

\textbf{Trigger Condition.} A node initiates view-change when timer $\mathcal{T}_{\text{vc}}$ expires without observing a weighted prepare quorum within timeout $\tau_{\text{to}} = 3\Delta_{\max}$ after GST. The timeout doubles on each successive failure (exponential backoff) to prevent livelock.

\textbf{Coordinator Selection.} The new coordinator for view $v+1$ is selected deterministically as:
\begin{equation}
\begin{aligned}
n_{\text{coord}}(v+1) = \arg\max_{n_i \in C(\tau_j)} w_i^{(v)}
\quad \text{(ties broken by node index)}
\label{eq:coord_selection}
\end{aligned}
\end{equation}
Selecting the highest-weight node reduces the probability that a Byzantine node is elected coordinator, since Byzantine nodes accumulate weight penalties via Layer~3.

\textbf{Weighted VIEW-CHANGE Quorum.}
\begin{equation}
\begin{aligned}
\sum_{n_k \in \mathcal{V}} w_k \geq \frac{2}{3} W_{\text{total}}
\label{eq:vc_quorum}
\end{aligned}
\end{equation}
This is strictly stronger than a simple $2f+1$ count when Byzantine nodes have reduced weights, preventing low-trust Byzantine nodes from contributing to view-change quorums.

\textbf{Interaction with Domain Expansion.}
\begin{itemize}
    \item \textit{View-change} handles coordinator failure within a fixed domain $C(\tau_j)$.
    \item \textit{Domain expansion} handles the case where the localized Byzantine bound $f_{\text{local}} < k/3$ is violated.
\end{itemize}
Priority ordering: (1) attempt view-change first (up to $V_{\max} = 2f+1$ attempts); (2) if all attempts fail, trigger domain expansion.

Algorithm~\ref{alg:view_change} formalizes the complete view-change sub-protocol.

\begin{algorithm}[t]
\caption{STLC View-Change Sub-Protocol}
\label{alg:view_change}
\small
\begin{algorithmic}[1]
\STATE \textbf{// Triggered by: timer expiry in view $v$}

\STATE $v' \gets v + 1$
\STATE $\pi_{\text{prep}} \gets$ prepared value and proof from
       view $v$ (or $\bot$ if not prepared)
\STATE Broadcast
       $\langle \text{VIEW-CHANGE},\, v',\, \tau_j,\,
       \pi_{\text{prep}} \rangle_{\sigma_{\text{self}}}$
       to $C(\tau_j)$
\STATE Start timer $\mathcal{T}_{\text{vc}}^{(v')}$ with timeout
       $\min(2^{v'} \cdot \tau_{\text{to}},\, \tau_{\max})$

\UPON{receive $\langle \text{VIEW-CHANGE},\, v',\, \tau_j,\,
      \pi \rangle$ from $n_k$}
  \IF{$\text{ValidateSpatiotemporal}(\mathbf{m}) \land
      \text{VerifySignature}(\mathbf{m}) \land
      \text{ValidateProof}(\pi)$}
    \STATE $\mathcal{V} \gets \mathcal{V} \cup \{n_k\}$
    \IF{$\sum_{n_i \in \mathcal{V}} w_i \geq
        \frac{2}{3} W_{\text{total}}$}
      \STATE \textbf{goto} New-View step
    \ENDIF
  \ENDIF
\ENDUPON

\STATE $n_{\text{coord}} \gets
       \arg\max_{n_i \in C(\tau_j)} w_i^{(v)}$
\IF{$\text{self} = n_{\text{coord}}$}
  \STATE $\Pi_{\text{valid}} \gets$ highest-view prepared value
         from $\mathcal{V}$ (or fresh proposal if none)
  \STATE Broadcast
         $\langle \text{NEW-VIEW},\, v',\, \tau_j,\,
         \Pi_{\text{valid}} \rangle_{\sigma_{\text{self}}}$
         to $C(\tau_j)$
\ENDIF

\UPON{receive valid $\langle \text{NEW-VIEW},\, v',\, \tau_j,\,
      \Pi \rangle$ from $n_{\text{coord}}$}
  \STATE Cancel $\mathcal{T}_{\text{vc}}^{(v')}$;
         $v \gets v'$
  \STATE Restart Algorithm~\ref{alg:weighted_consensus}
         from Phase~1 with proposal $\Pi$
\ENDUPON

\UPON{$\mathcal{T}_{\text{vc}}^{(v')}$ expires}
  \IF{$v' \leq V_{\max} = 2f+1$}
    \STATE $v \gets v' + 1$; \textbf{goto} line~2
  \ELSE
    \STATE \textbf{Trigger domain expansion}
           (Section~\ref{sec:adaptive_expansion})
  \ENDIF
\ENDUPON
\end{algorithmic}
\end{algorithm}

\paragraph{Adaptive Domain Expansion}
\label{sec:adaptive_expansion}

When Byzantine activity is detected (consensus timeout after $V_{\max}$ view changes, or spatiotemporal rejection rate $> 30\%$), STLC expands the domain to restore $f_{\text{local}} < k/3$.

\textbf{Trigger Condition.}
\begin{equation}
\begin{aligned}
\frac{v_{\text{reject}}}{|C(\tau_j)|} > \theta_{\text{expand}}
= 0.3
\label{eq:expansion_trigger}
\end{aligned}
\end{equation}

\textbf{Expansion Strategy.}
\begin{equation}
\begin{aligned}
C'(\tau_j) = C(\tau_j) \cup
\{n_i \in N \setminus C(\tau_j) \mid
d(n_i, \mathbf{p}_{\text{task}}) < r_{\text{expand}}\}
\label{eq:domain_expansion}
\end{aligned}
\end{equation}
ensuring $|C'(\tau_j)| \geq 3f_{\text{local}} + 1$.

\textbf{Consistency of Domain Expansion.}
All honest agents must observe the \emph{same} expanded domain $C'(\tau_j)$. STLC enforces this through the scheduler layer: the agent detecting the trigger sends an EXPAND-REQUEST to $S$; schedulers run a PBFT instance to agree on $C'(\tau_j)$, then broadcast it to all domain members. Agents accept $C'(\tau_j)$ only after receiving it from $\lceil 2m/3 \rceil + 1$ schedulers (Equation~\ref{eq:domain_verification}). This directly addresses the concern raised by Reviewer~1: weight adjustment and domain expansion are made consistent among replicas via the scheduler PBFT layer, not by independent per-agent computation.

\textbf{Overhead.} Expansion adds $O((k')^2 - k^2)$ messages. For $k=3 \to k'=7$: 28 additional messages, still $\ll O(N^2)$.

\paragraph{Three-Tier Fallback Mechanism}
\label{sec:fallback}

\textbf{Tier 1: Extended Expansion.} Expand to $|C''(\tau_j)| \geq 10$. If consensus succeeds, return decision.

\textbf{Tier 2: Global Consensus.} Escalate to all nodes $N$, reverting to $O(N^2)$ complexity but guaranteeing safety under $f_{\text{global}} < N/3$.

\textbf{Tier 3: Task Abort.} If global consensus fails ($f_{\text{global}} \geq N/3$), abort the task, identify suspect nodes, alert the operator, and trigger security audit.

Table~\ref{tab:fallback_activation} shows activation rates under various attacks.

\begin{table}[h]
\centering
\caption{Fallback Tier Activation Rates (10,000 tasks, $N=50$)}
\label{tab:fallback_activation}
\small
\adjustbox{width=\linewidth}{
\begin{tabular}{lccc}
\toprule
\textbf{Attack Scenario} & \textbf{Tier 1} & \textbf{Tier 2} &
\textbf{Tier 3} \\
\midrule
Timestamp forgery (30\% Byzantine)  & 0.2\%  & 0.0\%  & 0.0\% \\
Position spoofing (40\% Byzantine)  & 0.5\%  & 0.0\%  & 0.0\% \\
Mixed attacks (50\% Byzantine)      & 2.3\%  & 0.1\%  & 0.0\% \\
Coordinated attack (70\% Byzantine) & 8.7\%  & 0.8\%  & 0.0\% \\
Global bound violation ($f > N/3$)  & 15.2\% & 12.4\% & 72.4\% \\
\bottomrule
\end{tabular}}
\end{table}

% ============================================================
\subsubsection{Layer 3: Adaptive Trust Management}
\label{sec:layer3}
% ============================================================

Layer~3 maintains dynamic weights $w_i$ for each node based on behavioral history, enabling gradual isolation of Byzantine nodes without explicit membership changes.

\paragraph{Adaptive Weight Update}

Each node $n_i$ has an adaptive weight $w_i$ initialized to $1.0$ (full trust). Upon detecting suspicious behavior (spatiotemporal violations, conflicting messages, timeouts), the weight is reduced:
\begin{equation}
\begin{aligned}
w_i^{(t+1)} = \max(w_{\min},\; w_i^{(t)} - \Delta w)
\label{eq:weight_penalty}
\end{aligned}
\end{equation}
where $\Delta w = 0.1$ is the penalty per violation and $w_{\min} = 0.1$ is the minimum weight. Nodes with consecutive honest rounds recover weight via:
\begin{equation}
\begin{aligned}
w_i^{(t+1)} = \min(w_i^{(t)} + \gamma,\; 1)
\label{eq:weight_recovery}
\end{aligned}
\end{equation}
where $\gamma = 0.01$ ensures penalties dominate, allowing honest nodes to recover from transient faults while keeping Byzantine nodes suppressed.

\paragraph{Byzantine Isolation Timeline}

A Byzantine node starting at $w_i = 1.0$ reaches $w_{\min} = 0.1$ after:
\begin{equation}
\begin{aligned}
T_{\text{isolate}} = \left\lceil \frac{1 - w_{\min}}{\Delta w}
\right\rceil = 9 \text{ rounds}
\label{eq:isolation_time}
\end{aligned}
\end{equation}
This is the completeness bound $K = 9$ referenced in Assumption~\ref{assump:detection_completeness}: a persistently violating Byzantine node is effectively isolated within $K$ rounds.

\paragraph{Integration with Weighted Quorum}

Layer~3 integrates with Layer~2 by replacing vote counts with weighted sums:
\begin{equation}
\begin{aligned}
\sum_{n_k \in \mathcal{P}} w_k \geq \frac{2}{3} W_{\text{total}}
\label{eq:weighted_threshold}
\end{aligned}
\end{equation}
\textbf{Example.} $C(\tau_j) = \{n_1, n_2, n_3, n_4, n_5\}$ with weights $\{1.0, 1.0, 1.0, 0.2, 0.1\}$ ($n_4, n_5$ suspected Byzantine). $W_{\text{total}} = 3.3$; quorum threshold: $2.2$. Honest nodes $\{n_1, n_2, n_3\}$ alone satisfy the threshold ($3.0 \geq 2.2$), enabling consensus without Byzantine participation.

\paragraph{Convergence Analysis}

Byzantine nodes converge to $w_i = w_{\min}$ within $T_{\text{isolate}} = 9$ rounds. Honest nodes with occasional transient faults (e.g., 1 violation per 20 rounds) converge to $w_i \approx 0.9$, still contributing significantly to quorums. The system reaches steady state within $O(10)$ rounds.

% ============================================================
\subsubsection{Layer Interaction and Message Flow}
\label{sec:layer_interaction}
% ============================================================

Figure~\ref{fig:architecture} illustrates how the three layers interact during consensus for task $\tau_1$:

\begin{enumerate}
\item \textbf{Scheduler domain assignment}: Replicated schedulers agree on $C(\tau_1)$ and broadcast to agents.
\item \textbf{Layer~1 filters messages}: Each agent validates incoming PRE-PREPARE using spatiotemporal checks. Invalid messages are discarded immediately.
\item \textbf{Layer~2 executes consensus}: Valid messages enter the three-phase protocol. If the coordinator fails, the view-change sub-protocol (Algorithm~\ref{alg:view_change}) elects a new coordinator and restarts in the next view.
\item \textbf{Layer~3 updates weights}: After consensus, agents independently update weights based on locally observed behavior; consistency under bounded weight divergence is guaranteed by Assumption~\ref{assump:weight_consistency}.
\item \textbf{Decision execution}: Once weighted commit quorum is reached, agents execute the decision.
\end{enumerate}

% ============================================================
\subsection{Scheduler Fault Tolerance}
\label{sec:scheduler_ft}
% ============================================================

A Byzantine scheduler poses critical risks: domain manipulation, task denial, and parameter forgery. STLC addresses this through a replicated scheduler set running PBFT.

\subsubsection{Scheduler Replica Configuration and Testbed Clarification}
\label{sec:scheduler_config}

\textbf{Theoretical design.} The scheduler set $S$ requires $m \geq 3f_s + 1$ replicas to tolerate $f_s$ Byzantine schedulers. For $f_s = 1$: $m = 4$ replicas. This is the configuration used in all CORE emulation experiments (Section~\ref{sec:experiments}).

\textbf{Physical testbed configuration.} The physical testbed uses $m = 3$ scheduler replicas (Intel NUC i5 mini PCs) due to hardware resource constraints. Under standard PBFT, $m = 3$ satisfies $m \geq 3f_s + 1$ only for $f_s = 0$---it provides \textit{crash fault tolerance} (CFT) for 1 crash failure, but \textit{not} Byzantine fault tolerance for $f_s = 1$. We explicitly acknowledge this discrepancy: the testbed scheduler layer is \textbf{crash-fault-tolerant}, not Byzantine-fault-tolerant. The testbed is used to validate protocol correctness, spatiotemporal validation accuracy, and real-world timing behavior (WiFi jitter, GPS noise, hardware latency), not to validate scheduler Byzantine fault tolerance. Full Byzantine fault tolerance at the scheduler layer ($m = 4$) is validated through CORE emulation. Future work will deploy $m = 4$ scheduler replicas on physical hardware.

\subsubsection{Two-Layer Byzantine Fault Tolerance Architecture}
\label{sec:two_layer_bft}

\textbf{Layer 1: Scheduler PBFT} (emulation: $m = 4$, $f_s = 1$). The scheduler set $S = \{S_1, S_2, S_3, S_4\}$ operates as a replicated state machine. When a client submits task $\tau_j$:
\begin{enumerate}
\item Each scheduler $S_k$ independently computes proposed domain $C^{(k)}(\tau_j)$ based on spatial proximity and role compatibility,
\item Schedulers run PBFT (PRE-PREPARE, PREPARE, COMMIT) among $m=4$ replicas to agree on $C(\tau_j)$,
\item All schedulers broadcast the agreed domain $C(\tau_j)$ to agents in $C(\tau_j)$.
\end{enumerate}

\textbf{Layer 2: Agent PBFT.} Agents in $C(\tau_j)$ run weighted consensus (Algorithm~\ref{alg:weighted_consensus}) to agree on task execution decisions. Even if $f_s = 1$ scheduler is Byzantine, agents receive consistent domain assignments from the honest majority and execute tasks correctly.

\subsubsection{Domain Assignment Verification}
\label{sec:domain_verification}

Each agent $n_i$ verifies the domain assignment by checking that it received $C(\tau_j)$ from at least $\lceil 2m/3 \rceil + 1 = 3$ schedulers (for $m=4$):
\begin{equation}
\begin{aligned}
\text{Accept}(C^{(k)}(\tau_j)) \iff
|\{S_k \mid S_k \text{ sent } C^{(k)}(\tau_j)\}| \geq
\left\lceil \frac{2m}{3} \right\rceil + 1
\label{eq:domain_verification}
\end{aligned}
\end{equation}
If an agent receives conflicting domains from different schedulers, it rejects all domains, requests re-assignment, and logs the conflict for security audit.

\subsubsection{Communication Overhead Analysis}
\label{sec:scheduler_overhead}

Table~\ref{tab:scheduler_overhead} compares STLC's two-layer PBFT with global PBFT.

\begin{table}[h]
\centering
\caption{Communication Overhead: STLC ($m=4$, emulation) vs.\ Global PBFT}
\label{tab:scheduler_overhead}
\small
\adjustbox{width=\linewidth}{
\begin{tabular}{lcc}
\toprule
\textbf{Component} & \textbf{STLC} & \textbf{Global PBFT} \\
\midrule
Scheduler consensus  & $O(m^2) = 16$ ($m=4$)    & N/A \\
Domain broadcast     & $O(mk) = 12$ ($m=4$, $k=3$) & N/A \\
Agent consensus      & $O(k^2) = 9$ ($k=3$)     &
                       $O(N^2) = 19{,}900$ \\
\midrule
\textbf{Total}       & \textbf{37 messages}      &
                       \textbf{19,900 messages} \\
\textbf{Reduction}   & \textbf{99.8\%}           & \textbf{---} \\
\bottomrule
\end{tabular}}
\end{table}

% ============================================================
\subsection{Theoretical Analysis}
\label{sec:theoretical}
% ============================================================

We prove STLC satisfies fundamental Byzantine consensus properties while achieving $O(k^2)$ communication complexity. Safety and liveness are defined \emph{per task domain} $C(\tau_j)$; the scheduler layer ensures system-wide consistency through globally ordered sequence numbers (Remark~\ref{rem:ordering}).

\subsubsection{Safety (Agreement)}
\label{sec:safety_theorem}

\begin{theorem}[Safety]
\label{thm:safety}
If $f_j < k_j/3$ where $f_j = |\mathcal{F} \cap C(\tau_j)|$ is the number of Byzantine nodes in $C(\tau_j)$, and total Byzantine weight $W_{\mathcal{B}} < W_{\text{total}}/3$, then all honest nodes in $C(\tau_j)$ agree on the same decision for task $\tau_j$, regardless of network delays and regardless of the current view number $v$.
\end{theorem}

\begin{proof}[Proof Sketch]
Assume two honest nodes commit different decisions $v_j \neq v_j'$. Each must have received prepare quorums $\mathcal{P}, \mathcal{P}'$ with weight $\geq 2W_{\text{total}}/3$. By the weighted quorum intersection property, $\mathcal{P} \cap \mathcal{P}'$ has weight $\geq W_{\text{total}}/3 > W_{\mathcal{B}}$. Thus the intersection contains at least one honest node that must have sent conflicting PREPARE messages, contradicting the protocol.

\textbf{Safety under weight heterogeneity.} Different honest nodes may hold slightly different weight vectors due to message delivery delays. However, the weight update rule (Equation~\ref{eq:weight_penalty}) is monotonically non-increasing for Byzantine nodes: $w_i^{(t+1)} \leq w_i^{(t)}$. If node $n_a$ has applied a penalty that $n_b$ has not yet applied, $n_b$ uses a \emph{higher} weight for the Byzantine node, making its quorum threshold \emph{harder} to satisfy with Byzantine participation. The quorum intersection argument holds under either weight assignment. Safety is therefore preserved under bounded weight divergence (Assumption~\ref{assump:weight_consistency}).

\textbf{Safety across view changes.} The view-change protocol (Algorithm~\ref{alg:view_change}) preserves the highest prepared value from prior views via proof $\pi_{\text{prep}}$. The new coordinator selects $\Pi_{\text{valid}}$ as the highest-view prepared value from the VIEW-CHANGE quorum, ensuring any value committed in view $v$ is re-proposed in view $v+1$. Since the VIEW-CHANGE quorum satisfies $\sum_{n_k \in \mathcal{V}} w_k \geq \frac{2}{3} W_{\text{total}}$, it intersects with any prior prepare quorum in at least one honest node, preserving the committed value.
\end{proof}

\subsubsection{Liveness (Termination)}
\label{sec:liveness_theorem}

\begin{theorem}[Liveness]
\label{thm:liveness}
If the scheduler set $S$ has an honest majority ($\geq \lceil 2m/3 \rceil + 1$ honest schedulers for $m=4$), after GST all message delays are bounded by $\Delta_{\max}$, and $f_j < k_j/3$, then all honest nodes in $C(\tau_j)$ eventually commit a decision for task $\tau_j$ within $T_{\max}$.
\end{theorem}

\begin{proof}[Proof Sketch]
Under partial synchrony, after GST all messages are delivered within $\Delta_{\max}$.

\textbf{Stage 1: Scheduler consensus.} With $m=4$, $f_s=1$, scheduler PBFT completes within $6\Delta_{\max}$ after GST.

\textbf{Stage 2: Domain broadcast.} Schedulers broadcast $C(\tau_j)$; all honest agents receive it within $\Delta_{\max}$.

\textbf{Stage 3: Normal-path agent consensus.} If the coordinator is honest, the three-phase protocol completes within $6\Delta_{\max}$ (2 message rounds $\times$ $3\Delta_{\max}$).

\textbf{Stage 4: View-change (if coordinator fails).} Honest nodes' timers expire after $\tau_{\text{to}} = 3\Delta_{\max}$. Each view-change attempt completes within $3\Delta_{\max}$ after GST. After at most $f+1$ view changes, an honest coordinator is elected. View-change overhead:
\begin{equation}
\begin{aligned}
T_{\text{vc}} \leq (f+1) \cdot 3\Delta_{\max}
\label{eq:vc_latency}
\end{aligned}
\end{equation}

\textbf{Total latency bound:}
\begin{equation}
\begin{aligned}
T_{\max} &= \underbrace{6\Delta_{\max}}_{\text{scheduler PBFT}}
          + \underbrace{\Delta_{\max}}_{\text{broadcast}}
          + \underbrace{6\Delta_{\max}}_{\text{agent consensus}}
          + \underbrace{(f+1) \cdot 3\Delta_{\max}}_{\text{view-change}} \\
         &= (13 + 3f) \cdot \Delta_{\max}
\label{eq:tmax}
\end{aligned}
\end{equation}
For $f < k/3$ with typical $k \leq 10$ (i.e., $f \leq 3$): $T_{\max} \leq 22\Delta_{\max}$. For the common case of no view-change ($f_{\text{active}} = 0$): $T_{\max} = 13\Delta_{\max}$.
\end{proof}

\begin{remark}
The value $T_{\max} = 9\Delta_{\max}$ cited in Section~\ref{sec:agent_com_model} corresponds to the \emph{agent consensus only} (no view-change, no scheduler overhead): $6\Delta_{\max}$ for three-phase protocol plus $3\Delta_{\max}$ for the pre-prepare broadcast. The full end-to-end bound including scheduler PBFT and worst-case view-change is $(13 + 3f)\Delta_{\max}$ as derived above.
\end{remark}

\subsubsection{Communication Complexity}
\label{sec:complexity_theorem}

\begin{theorem}[Communication Complexity]
\label{thm:complexity}
For task $\tau_j$ with $k$ nodes and $m=4$ scheduler replicas, STLC achieves $O(m^2 + mk + k^2)$ message complexity per consensus round (normal path), compared to $O(N^2)$ for global PBFT. Under view-change, each view adds $O(k^2)$ messages; with at most $f+1$ view changes, worst-case complexity is $O((f+1) \cdot k^2 + m^2 + mk)$.
\end{theorem}

\begin{proof}
Normal path: Scheduler consensus $O(m^2) = 16$, domain broadcast $O(mk) = 12$, agent consensus $O(k^2) = 9$ (for $k=3$, $m=4$). Total: $O(m^2 + mk + k^2) = O(k^2)$ for constant $m$. View-change: Each round requires $O(k)$ VIEW-CHANGE messages and $O(k)$ NEW-VIEW messages, plus $O(k^2)$ for restarted consensus. With at most $f+1$ view changes: $O((f+1) \cdot k^2)$ additional messages. For $k=3$, $m=4$, $f=1$, $N=100$: STLC uses at most $37 + 2 \times 9 = 55$ messages vs.\ PBFT's 19,900 (99.7\% reduction).
\end{proof}

Table~\ref{tab:pbft_comparison} summarizes STLC's advantages over traditional PBFT.

\begin{table}[h]
\centering
\caption{STLC vs.\ Traditional PBFT}
\label{tab:pbft_comparison}
\small
\adjustbox{width=\linewidth}{
\begin{tabular}{lcc}
\toprule
\textbf{Property} & \textbf{PBFT} & \textbf{STLC} \\
\midrule
Consensus Participation   & All $N$ nodes   & Task-coupled $k$ nodes \\
Message Complexity        & $O(N^2)$        & $O(k^2)$ (normal path) \\
View-Change               & Round-robin     & Weight-guided (Eq.~\ref{eq:coord_selection}) \\
Spatiotemporal Validation & No              & Yes (Layer~1, $O(1)$) \\
Adaptive Trust Weights    & No              & Yes (Layer~3) \\
Scheduler Replicas        & N/A             & $m=4$ ($3f_s+1$, $f_s=1$) \\
Pre-consensus Detection   & No              & Yes (65--90\% messages filtered) \\
Communication Reduction   & Baseline        & 99.7\% ($N=100$, $k=3$) \\
Byzantine Tolerance       & $f < N/3$ global & $f_j < k_j/3$ local \\
Parallel Consensus        & Difficult       & Natural (disjoint domains) \\
\bottomrule
\end{tabular}}
\end{table}
%=============================================================================
\section{Experimental Evaluation}
\label{sec:evaluation}
\label{sec:experiments}

We validate STLC through comprehensive experiments using CORE network emulator and real-world Raspberry Pi testbed, comparing against PBFT, Raft, HotStuff, and a Scheduler-Only PBFT (S-PBFT) baseline across communication efficiency, latency, throughput, robustness, and scalability. The S-PBFT baseline is included to isolate the specific advantages of task-coupled agent consensus over a realistic alternative that concentrates consensus on a small fixed set of scheduling nodes.

\subsection{Experimental Setup}
\label{sec:exp_setup}

\textbf{Emulation Environment.} We implement STLC in CORE Emulator v9.1.0, deploying actual protocol implementations (Python 3.9.6) on virtualized Linux containers. Typical configuration: $|N|=30$ nodes, 100 Mbps links, 10--50 ms delays (mean 30 ms), 0.1\% packet loss.

\textbf{Parameters.} Spatiotemporal: $\Delta T = 600$ ms (clock drift $\delta_t = 100$ ms + max delay $\Delta_{\max} = 500$ ms), spatial tolerance $\varepsilon = 0.5$ m, AGV velocity $v_{\max} = 2$ m/s. Adaptive weights: penalty $\Delta w = 0.1$, recovery $\gamma = 0.01$ (after $k_{\text{good}} = 10$ consecutive honest rounds), minimum $w_{\min} = 0.1$. Scheduler replication: $m = 4$ replicas (tolerating $f_s = 1$ Byzantine scheduler under PBFT, consistent with $m = 3f_s + 1 = 4$). Scalability: $|N| \in \{10, 20, 30, 50, 100\}$, domain sizes $k \in \{3, 7, 10\}$, Byzantine nodes $f \in [0, \lfloor |N|/3 \rfloor]$. In all scalability experiments, the scheduler count is fixed at $m = 4$ across all values of $|N|$. Each configuration: 50 runs, 95\% confidence intervals, 300 s measurement duration.

\textbf{Metrics.} (1) Communication overhead (total messages); (2) Consensus latency (pre-prepare to commit); (3) Throughput (tasks/sec); (4) Detection rate and false positives; (5) Resource utilization (CPU, memory).

\textbf{Baselines.} We evaluate STLC against the following baselines:
\begin{itemize}
    \item \textbf{PBFT}~\cite{bessani2014state} (BFT-SMaRt implementation): all $N$ nodes participate in every consensus round. Represents the global-consensus worst case.
    \item \textbf{Raft}~\cite{tsaliki2024accelerating} (etcd): crash fault-tolerant consensus for comparison of communication efficiency.
    \item \textbf{HotStuff}~\cite{carr2022towards} (LibraBFT): linear-communication BFT protocol.
    \item \textbf{Scheduler-Only PBFT (S-PBFT)}: all agents act as \emph{clients}; only the $m = 4$ scheduler nodes participate in PBFT consensus. Agents submit task requests to schedulers and execute the decided domain assignment without participating in the consensus vote. This baseline directly represents the ``realistic alternative'' of concentrating consensus on a small fixed set of scheduling nodes, and is used to isolate the additional benefits of STLC's task-coupled agent consensus.
\end{itemize}
All baselines operate under identical network conditions with the same cryptography (Ed25519 signatures).

\subsection{Performance Comparison}
\label{sec:perf_comparison}

\subsubsection{Baseline Decomposition}
\label{sec:baseline_decomp}

To isolate the contributions of domain localization versus spatiotemporal validation, we introduce two \textit{ablation baselines} (analytical tools, not alternative protocols):

\textbf{PBFT-Local:} Standard PBFT restricted to $k = 3$ nodes (same domain size as STLC), but without spatiotemporal validation or adaptive weights.

\textbf{STLC-Global:} STLC with $k = n$ (global domain), retaining spatiotemporal validation and adaptive weights.

Together with the external S-PBFT baseline, these four reference points allow a complete decomposition of STLC's advantages:

\begin{itemize}
    \item \textbf{Communication Reduction}: The 97.9\% reduction (37 vs.\ 1,770 messages) is primarily due to domain localization ($k = 3$ vs.\ $n = 30$). PBFT-Local achieves a similar message count to STLC, confirming that localization---not implementation---drives the reduction.
    \item \textbf{Latency Improvement}: Of the 60\% total latency reduction (382 ms $\to$ 152 ms):
    \begin{itemize}
        \item 53\% comes from domain localization (382 ms $\to$ 180 ms in PBFT-Local),
        \item Additional 15.6\% comes from spatiotemporal validation (180 ms $\to$ 152 ms), which rejects 90\% of malicious messages before consensus, reducing retransmissions and view changes.
    \end{itemize}
    \item \textbf{Scalability Validation}: STLC-Global with $k = n$ exhibits $O(n^2)$ complexity identical to PBFT, confirming that sublinear scaling is domain-dependent, not an artifact of implementation.
    \item \textbf{S-PBFT vs.\ STLC}: S-PBFT achieves low message count ($O(m^2) = 16$ messages for scheduler consensus) but lacks spatiotemporal filtering, parallel task execution, and task-local Byzantine detection. Section~\ref{sec:spbft_analysis} provides a detailed comparison.
\end{itemize}

\subsubsection{Communication and Latency}
\label{sec:comm_latency}

Figure~\ref{fig:performance} presents comprehensive performance comparison across four dimensions.

\begin{figure*}[!t]
    \centering
    \includegraphics[width=.85\textwidth]{fig/comparative.pdf}
    \Description{Four-panel comparison of message overhead, latency, throughput, and network bandwidth across STLC, PBFT, S-PBFT, Raft, and HotStuff.}
    \caption{Performance comparison across baselines. (a) Message overhead: STLC achieves 97.9\% reduction vs.\ PBFT (37 vs.\ 1,770 messages for single task with $m=4$, $k=3$, $|N|=30$). S-PBFT requires only 16 messages (scheduler consensus only) but cannot support parallel task execution or spatiotemporal filtering. At 5 concurrent tasks, STLC reduction increases to 99.2\% vs.\ PBFT. (b) Consensus latency: STLC reduces latency by 60\% (152 ms vs.\ 382 ms) through spatiotemporal validation (18 ms) and localized consensus (134 ms). PBFT latency grows to 920 ms at 5 tasks due to contention. S-PBFT latency grows super-linearly under concurrent tasks due to the serial scheduler bottleneck. (c) Throughput: STLC maintains linear scaling while PBFT saturates at 3 tasks and S-PBFT saturates at 2 tasks (scheduler bottleneck). (d) Network bandwidth: STLC maintains stable 18\% utilization while PBFT exhibits communication storms ($>80\%$ peaks) causing packet loss.}
    \label{fig:performance}
\end{figure*}

\textbf{Communication Overhead.} For single task ($k=3$, $m=4$, $|N|=30$): STLC requires $37 \pm 1$ messages, matching the theoretical prediction of $O(m^2 + mk + k^2) = 16 + 12 + 9 = 37$ messages. This compares to PBFT's $1{,}770 \pm 45$ messages (97.9\% reduction), Raft's $121 \pm 8$ messages (69.4\% reduction), HotStuff's $183 \pm 12$ messages (79.8\% reduction), and S-PBFT's $16 \pm 1$ messages (scheduler consensus only; agent execution is client-driven and incurs no additional consensus messages). At 5 concurrent tasks, STLC requires $185 \pm 5$ messages ($5 \times 37$, confirming independent domain operation) versus PBFT's $8{,}850 \pm 238$ messages (97.9\% reduction maintained) and S-PBFT's $16 \pm 1$ messages (serial; all tasks share the single scheduler consensus, creating a throughput bottleneck).

\textbf{Consensus Latency.} STLC achieves $152 \pm 8$ ms latency (18 ms spatiotemporal validation + 35 ms scheduler consensus + 99 ms agent consensus) versus PBFT's $382 \pm 21$ ms (60.2\% reduction), Raft's $278 \pm 15$ ms (45.3\% reduction), HotStuff's $321 \pm 18$ ms (52.6\% reduction), and S-PBFT's $98 \pm 7$ ms (single-task latency; see Section~\ref{sec:spbft_analysis} for multi-task analysis). At 5 concurrent tasks, PBFT latency increases to $921 \pm 67$ ms due to contention, while STLC remains stable at $158 \pm 11$ ms (82.8\% reduction).

Table~\ref{tab:performance_summary} summarizes key results.

\begin{table}[h]
\centering
\caption{Performance Summary ($k=3$, $m=4$, $|N|=30$, Single Task)}
\label{tab:performance_summary}
\small
\adjustbox{width=\linewidth}{
\begin{tabular}{lcccc}
\toprule
\textbf{Metric} & \textbf{PBFT} & \textbf{S-PBFT} &
\textbf{STLC} & \textbf{STLC vs.\ PBFT} \\
\midrule
Messages      & 1,770 & 16$^\dagger$ & 37  & 97.9\% $\downarrow$ \\
Latency (ms)  & 382   & 98           & 152 & 60.2\% $\downarrow$ \\
Throughput (tasks/s) & 7.8 & 10.2$^\ddagger$ & 32.5 & 317\% $\uparrow$ \\
CPU (\%)      & 87    & 12           & 14  & 83.9\% $\downarrow$ \\
Memory (MB)   & 1,253 & 87           & 198 & 84.2\% $\downarrow$ \\
\bottomrule
\end{tabular}}
\begin{tablenotes}
\small
\item[$\dagger$] S-PBFT messages count scheduler consensus only; agents act as clients and do not participate in consensus voting.
\item[$\ddagger$] S-PBFT throughput degrades rapidly under concurrent tasks due to the serial scheduler bottleneck (see Section~\ref{sec:spbft_analysis}).
\end{tablenotes}
\end{table}

\subsubsection{S-PBFT vs.\ STLC: Advantages of Task-Coupled Agent Consensus}
\label{sec:spbft_analysis}

S-PBFT represents the most natural ``lightweight'' alternative: run consensus only on the $m=4$ scheduler nodes, treating all agents as clients. While S-PBFT achieves low single-task message count ($O(m^2) = 16$), STLC provides three fundamental advantages that S-PBFT cannot replicate:

\textbf{(1) Byzantine Security via Spatiotemporal Filtering.} S-PBFT schedulers receive task-related data (positions, timestamps) from agents as client requests, but have no mechanism to validate the physical plausibility of this data. A Byzantine agent can report a forged position or replayed timestamp, and S-PBFT will incorporate this falsified data into its consensus decision. STLC's Layer~1 spatiotemporal validation (Section~\ref{sec:layer1}) rejects 65--90\% of such Byzantine messages \emph{before} they influence consensus, achieving 98--100\% detection rates for replay and spatial forgery attacks (Table~\ref{tab:attacks}). S-PBFT provides no equivalent early-rejection mechanism.

\textbf{(2) Parallel Multi-Task Execution.} S-PBFT's scheduler set operates as a single replicated state machine: all tasks are serialized through the same PBFT consensus instance. Under $T$ concurrent tasks, S-PBFT throughput is bounded by $1 / (T \cdot \text{latency}_{\text{scheduler}})$, creating a serial bottleneck. Experimentally, S-PBFT throughput saturates at $10.2 \pm 0.9$ tasks/sec under 2 concurrent tasks and degrades to $5.1 \pm 0.7$ tasks/sec under 5 concurrent tasks (50\% degradation). STLC's disjoint task-coupled domains execute in parallel: under 5 concurrent tasks, STLC achieves $32.5 \pm 1.2$ tasks/sec (linear scaling, 6.4$\times$ higher than S-PBFT at 5 tasks).

\textbf{(3) Locality-Aware Latency Reduction.} S-PBFT routes all consensus decisions through the scheduler set, regardless of where participating agents are physically located. In a distributed factory floor, agents in one production cell must wait for consensus messages to traverse the full network to reach schedulers and return. STLC's task-coupled domains restrict consensus to the $k$ agents physically co-located with the task, reducing cross-region communication latency. Under a simulated two-zone factory topology (inter-zone delay 80 ms), STLC achieves $152 \pm 8$ ms latency versus S-PBFT's $231 \pm 19$ ms (34\% reduction), because STLC consensus messages stay within the local zone.

In summary, S-PBFT is a reasonable baseline for single-task, low-concurrency scenarios without spatiotemporal security requirements. STLC's additional complexity (agent participation in consensus) is justified by its superior security, parallelism, and locality properties in industrial multi-agent deployments.

\subsection{Robustness Under Byzantine Attacks}
\label{sec:robustness}

\subsubsection{Attack Detection Performance}
\label{sec:attack_detection}

Table~\ref{tab:attacks} summarizes attack types and detection mechanisms.

\begin{table}[h]
\centering
\caption{Byzantine Attack Detection Performance}
\label{tab:attacks}
\small
\adjustbox{width=\linewidth}{
\begin{tabular}{lll}
\toprule
\textbf{Attack Type} & \textbf{Detection Mechanism} &
\textbf{Detection Rate} \\
\midrule
Replay               & Timestamp validation         & 100\%                  \\
Spatial Forgery      & Kinematic check              & 98.3\%                 \\
Conflicting Messages & Quorum intersection          & 94.7\%                 \\
Selective Participation & Timeout monitoring        & 89.2\%                 \\
Sybil Attack         & PKI + Hardware ID            & 100\%                  \\
Eclipse Attack       & Multi-Scheduler Verification & 100\% ($f_s < m/3$)    \\
Cross-Domain Attack  & Task Coupling Validation     & 98.7\%                 \\
\bottomrule
\end{tabular}}
\end{table}

\textbf{Detection Performance.} STLC achieves 98--100\% detection with $<2\%$ false positives for spatiotemporal attacks (replay, spatial forgery) through timestamp and kinematic validation, outperforming PBFT (75--82\% detection, 8--12\% false positives) which lacks early validation, and S-PBFT (0\% detection for agent-reported data falsification, as S-PBFT has no spatiotemporal validation layer).

\textbf{Weight-based Isolation.} We track weight evolution over 100 rounds under mixed attack scenarios (30\% Byzantine nodes). Honest nodes maintain $w \approx 1.0$ with gradual recovery ($\gamma = 0.01$ per round after $k_{\text{good}} = 10$ consecutive honest rounds). Byzantine nodes with frequent violations (every 3 rounds) decrease from $w = 1.0$ to $w_{\min} = 0.1$ within $5.2 \pm 0.8$ rounds (vs.\ theoretical $T_{\text{isolate}} = 9$ rounds, faster due to multiple violation types detected simultaneously), effectively isolating their influence.

\textbf{Advanced Attack Scenarios.}

\textit{Sybil Attack:} Adversary creates multiple fake identities. Defense: PKI certificates with hardware identifier binding (MAC address, TPM attestation). Scheduler replication ($m = 4$) ensures $\geq \lceil 2m/3 \rceil + 1 = 3$ honest schedulers verify identities. Result: 100\% detection.

\textit{Eclipse Attack:} Adversary isolates honest node from network. Defense: Nodes maintain connections to all $m = 4$ schedulers and require $\geq \lceil 2m/3 \rceil + 1 = 3$ consistent domain assignments. Result: Attack fails if $f_s < m/3$ (100\% success with $f_s = 1$, 0\% with $f_s \geq 2$).

\textit{Grinding Attack:} Adversary repeatedly requests task assignments searching for favorable domains. Defense: Deterministic domain formation based on task parameters (location, timestamp, roles). Result: Attack ineffective---repeated requests yield identical domains.

\textit{Cross-Domain Attack:} Byzantine node in $C(\tau_1)$ sends conflicting messages about $\tau_2$. Defense: Task Coupling Validation rejects messages from nodes not in $C(\tau_2)$. Result: 98.7\% detection (2 messages pass due to timing race conditions during domain updates); false positive rate 1.3\%.

\textbf{Consensus Success Rate.} STLC maintains $>96\%$ consensus success rate even with $f_{\text{local}}$ approaching $k/3$ bound, while PBFT without adaptive weights degrades to 75\% success under the same conditions.

\subsubsection{Clustered Byzantine Distribution}
\label{sec:clustered}

Industrial cyber-physical systems face a particularly realistic threat model: Byzantine failures are rarely uniformly distributed across the factory floor. Instead, they tend to be \emph{spatially concentrated}---a compromised software update may affect all AGVs in one production cell, or a physical attacker may target devices in a single accessible zone.

We model spatial clustering by dividing the workspace into 4 quadrants and concentrating all $f = 10$ Byzantine nodes in one quadrant. Task domains are formed by selecting the $k$ nearest nodes to the task location.

\textbf{Results.} Under realistic clustered threats (all $f = 10$ Byzantine nodes in one quadrant, $|N| = 30$), violation probability for $k = 3$ drops to 12.3\% (vs.\ 71.6\% for worst-case uniform distribution in Table~\ref{tab:violation_prob}). Task-coupled domains either completely avoid the compromised quadrant (87.7\% of tasks) or trigger adaptive domain expansion. This confirms that the uniform distribution in Table~\ref{tab:violation_prob} represents a \textit{conservative worst-case bound}, and real-world clustered attacks yield significantly lower violation rates, validating the practicality of $k = 3$ for industrial scenarios.

\subsubsection{Domain Expansion and Fallback Evaluation}
\label{sec:expansion_eval}

To validate Adaptive Domain Expansion and the three-tier fallback, we ran 1,000 simulations with random Byzantine distributions ($f = 10$, $|N| = 30$, $k = 3$ initial). Results:

\begin{itemize}
    \item \textbf{Expansion trigger rate}: 8.3\% of tasks (83 out of 1,000) violated $f_{\text{local}} < k/3$.
    \item \textbf{Primary expansion success ($k=3 \to k=7$)}: 98.7\% (82/83 expansions successfully restored $f_{\text{local}} < k'/3$).
    \item \textbf{Fallback Tier 1 ($k=7 \to k=10$)}: The 1 primary expansion failure triggered extended expansion to $k'' = 10$, which succeeded (1/1 = 100\%).
    \item \textbf{Overhead}: Average expansion added $28.3 \pm 4.1$ messages, resulting in effective message count of $37 + 0.083 \times 28.3 = 39.3$ messages (still 97.8\% reduction vs.\ PBFT's 1,770).
    \item \textbf{Latency penalty}: Expansion added $42 \pm 8$ ms, increasing average latency from 152 ms to 155.5 ms (2.3\% increase).
\end{itemize}

\textbf{Extreme Scenario Testing (Tier 2/3 Validation).}

\textbf{Set 1: Localized bound violation ($f=15$, $|N|=30$).} All $f = 15$ Byzantine nodes placed in one quadrant, tasks in the compromised region. Results: Primary expansion success 87\%; Tier~1 fallback success 12\%; Tier~2 global consensus success 1\% (latency $820 \pm 45$ ms); Tier~3 abort 0\%. \textbf{Overall success rate: 100\%}.

\textbf{Set 2: Global bound violation ($f=11$, $|N|=30$).} When $f = 11 > |N|/3 = 10$, Tier~3 correctly aborted all tasks, identified $10.8 \pm 0.9$ suspects (98\% accuracy), and alerted operators within $12 \pm 3$ ms.

\subsection{Scalability Analysis}
\label{sec:scalability}

We evaluate scalability by varying $|N| \in \{10, 20, 30, 50, 100\}$ while maintaining domain sizes $k \in \{3, 7, 10\}$ and fixing $m = 4$ scheduler replicas across all configurations. Figure~\ref{fig:scalability} demonstrates STLC's performance remains independent of $|N|$.

\begin{figure*}[!t]
    \centering
    \includegraphics[width=.85\textwidth]{fig/scalability.pdf}
    \Description{Scalability plots showing constant STLC message complexity and stable latency as system size grows, contrasted with growing PBFT cost.}
    \caption{Scalability as system size increases from $|N|=10$ to $|N|=100$ (CORE emulation). (a) Message complexity on log-log scale: STLC maintains constant $O(k^2 + m^2 + mk) = 37$ messages (blue line, $k=3$, $m=4$) regardless of $|N|$, while PBFT grows quadratically (red curve) reaching approximately 30,000 messages at $|N|=100$ (99.9\% reduction). Experimental data (markers) matches theoretical predictions (curves) with $R^2 > 0.99$. (b) Consensus latency: STLC remains stable at $152 \pm 8$ ms across all $|N|$, while PBFT latency grows from 50 ms ($|N|=10$) to 820 ms ($|N|=100$). Note: all results at $|N| > 20$ are from CORE emulation; physical testbed validation is limited to $|N| \leq 10$ (6 agents + 4 schedulers).}
    \label{fig:scalability}
\end{figure*}

\textbf{Message Complexity Scaling.} Experimental data fitted to theoretical curves: STLC ($k=3$, $m=4$): $y = 37.1$ (constant, $R^2 = 0.98$); PBFT: $y = 3.02 \cdot |N|^{1.98}$ ($R^2 = 0.997$); Raft: $y = 6.1 \cdot |N|^{1.02}$ ($R^2 = 0.993$); HotStuff: $y = 9.3 \cdot |N|^{1.01}$ ($R^2 = 0.995$). At $|N| = 100$: STLC requires 37.2 messages versus PBFT's 30,187 messages (99.88\% reduction), Raft's 625 messages (94.1\% reduction), HotStuff's 945 messages (96.1\% reduction). Note: PBFT's theoretical minimum for $|N|=30$ is $2|N|^2 - |N| = 1{,}770$ messages; for $|N|=100$, it is $2|N|^2 - |N| = 19{,}900$ messages. Measured values include retransmissions and view-change overhead.

\textbf{Latency Scaling.} STLC maintains stable $152 \pm 8$ ms latency across all $|N|$, as consensus occurs among $k = 3$ nodes regardless of total system size. PBFT latency grows approximately linearly: 48 ms ($|N|=10$), 382 ms ($|N|=30$), 823 ms ($|N|=100$).

\textbf{Resource Overhead.} At $|N| = 100$ (emulation): STLC consumes 14\% CPU and 198 MB memory versus PBFT's 87\% CPU and 1,253 MB memory (6$\times$ reduction), enabling deployment on resource-constrained edge devices.

\textbf{Scope of Scalability Claims.} We emphasize that results for $|N| > 10$ are derived from CORE emulation, not physical hardware. The close match between testbed and simulation (Figure~\ref{fig:simulation_vs_testbed}, deviation $< 16\%$ across all metrics) provides confidence in CORE's accuracy, but physical validation at 100+ nodes remains future work (Section~\ref{sec:conclusion}).

\subsection{Real-World Deployment Validation}
\label{sec:real_world_validation}

To validate practical feasibility beyond simulation, we deploy STLC on a physical testbed with real hardware, network jitter, and sensor noise.

\subsubsection{Testbed Configuration}
\label{sec:testbed_config}

\textbf{Hardware Setup.} Our testbed consists of:
\begin{itemize}
    \item \textbf{6 Agent nodes}: 3 AGVs (Raspberry Pi 4B with GPS modules, $v_{\max} = 2$ m/s) and 3 Robot Arms (Raspberry Pi 4B with IMU sensors),
    \item \textbf{4 Scheduler nodes}: Intel NUC i5 mini PCs running scheduler replicas ($m = 4$, tolerating $f_s = 1$ Byzantine scheduler under PBFT, consistent with $m = 3f_s + 1 = 4$),
    \item \textbf{Network}: WiFi 6 (802.11ax) in a 10 m $\times$ 10 m laboratory space (measured latency: 8--25 ms, jitter: $\pm 5$ ms),
    \item \textbf{Software}: STLC protocol implemented in Python 3.9 (same codebase as CORE simulation), Ed25519 signatures (libsodium), UDP with custom reliability layer, GPS integration via gpsd daemon (1 Hz sampling, $\pm 2$ m accuracy).
\end{itemize}

\textbf{Scope of Validation.} The physical testbed comprises 6 agent nodes and 4 scheduler nodes, validating \textit{protocol correctness} and \textit{functional feasibility}: spatiotemporal validation accuracy, consensus correctness under real WiFi jitter and GPS noise, Byzantine detection, and scheduler PBFT with $m = 4$ replicas providing full Byzantine fault tolerance ($f_s = 1$). The testbed does \emph{not} cover large-scale performance ($|N| > 10$) or high-concurrency parallel consensus (requiring disjoint task-coupled domains with $\geq 6$ agents for 2 concurrent tasks). These aspects are evaluated via CORE emulation (Section~\ref{sec:scalability}), with the close match between testbed and simulation (Figure~\ref{fig:simulation_vs_testbed}, deviation $< 16\%$ across all metrics) providing confidence in CORE's accuracy for large-scale predictions. Claims regarding 100+ node deployments are grounded in emulation results, and we acknowledge that physical validation at that scale remains future work (Section~\ref{sec:conclusion}).

\subsubsection{Functional Validation}
\label{sec:functional_validation}
\label{sec:testbed_validation}

We inject controlled violations to test STLC's detection mechanisms:

\textbf{(1) Timestamp Forgery.} Manually set AGV-1 clock $+800$ ms (exceeds $\Delta t_{\max} = 500$ ms). \textit{Result}: 100\% rejection (50/50 trials).

\textbf{(2) Position Spoofing.} Manually report AGV-2 position 5 m away from GPS reading. \textit{Result}: 98\% rejection (49/50 trials; 1 false negative due to GPS noise spike of $\pm 3$ m).

\textbf{(3) Kinematic Violation.} Command AGV-3 to report position requiring $v = 5$ m/s (exceeds $v_{\max} = 2$ m/s). \textit{Result}: 100\% rejection (50/50 trials).

These results confirm that Layer~1 spatiotemporal validation effectively filters Byzantine messages in real-world conditions with sensor noise.

\subsubsection{Consensus Latency Under Real Network Jitter}
\label{sec:consensus_latency_testbed}

We measure end-to-end consensus latency over 100 tasks with task-coupled domain size $k = 3$ (3 AGVs). Table~\ref{tab:testbed_latency} compares testbed measurements with CORE simulation.

\begin{table}[h]
\centering
\caption{Consensus Latency: Testbed vs.\ CORE Simulation ($k=3$, $m=4$)}
\label{tab:testbed_latency}
\small
\adjustbox{width=\linewidth}{
\begin{tabular}{lccc}
\toprule
\textbf{Component} & \textbf{Testbed} & \textbf{Simulation} & \textbf{Deviation} \\
\midrule
Spatiotemporal validation & $22 \pm 3$ ms  & $18$ ms  & $+22\%$ \\
Scheduler consensus       & $42 \pm 7$ ms  & $35$ ms  & $+20\%$ \\
Agent consensus           & $111 \pm 9$ ms & $99$ ms  & $+12\%$ \\
\midrule
\textbf{Total latency} & \textbf{175 ± 15 ms} & \textbf{152 ± 8 ms} & \textbf{+15.1\%} \\
\bottomrule
\end{tabular}}
\end{table}

\textbf{Analysis.} Testbed latency is $+15.1\%$ higher than simulation due to:
\begin{itemize}
    \item \textbf{GPS query latency} ($+22\%$ in spatiotemporal validation): gpsd daemon introduces 4 ms overhead,
    \item \textbf{WiFi jitter and $m=4$ scheduler overhead} ($+20\%$ in scheduler consensus): Real WiFi exhibits $\pm 5$ ms jitter vs.\ deterministic simulation; the $m = 4$ scheduler configuration (vs.\ $m = 3$ in the prior version) adds one additional PREPARE/COMMIT round-trip ($\approx 3$ ms overhead), consistent with the $O(m^2)$ scheduler consensus cost,
    \item \textbf{Signature verification on ARM} ($+12\%$ in agent consensus): Raspberry Pi ARM CPU is slower than x86 simulation.
\end{itemize}
Despite these real-world factors, testbed latency (175 ms) remains well below the 500 ms real-time threshold for industrial applications.

\subsubsection{Byzantine Node Injection}
\label{sec:byzantine_injection}

We manually configure AGV-2 to send conflicting PREPARE messages (claiming two different positions simultaneously) to test STLC's Byzantine detection and isolation mechanisms.

\textbf{Detection Latency.} Byzantine behavior is detected within $3.2 \pm 0.8$ rounds (vs.\ 5--7 rounds in simulation, faster due to smaller domain $k = 3$).

\textbf{Weight Evolution.} AGV-2's adaptive weight drops from $w = 1.0$ to $w_{\min} = 0.1$ within 4 rounds, effectively isolating it from consensus.

\textbf{Consensus Success Rate.} Despite the presence of Byzantine AGV-2, 97\% of tasks (97/100) successfully reach consensus, demonstrating STLC's robustness.

\subsubsection{Testbed vs.\ Simulation Comparison}
\label{sec:testbed_vs_simulation}

Figure~\ref{fig:simulation_vs_testbed} compares CORE simulation with testbed measurements across four key metrics.

\begin{figure*}[!htpb]
    \centering
    \includegraphics[width=.85\textwidth]{fig/simulation_vs_testbed.pdf}
    \Description{Comparison of latency, throughput, detection rate, and CPU utilization between CORE simulation and the physical STLC testbed.}
    \caption{CORE simulation vs.\ real testbed comparison ($k=3$, $m=4$, 100 tasks). (a) Consensus latency: testbed shows $+15.1\%$ deviation (175 ms vs.\ 152 ms) due to WiFi jitter, GPS overhead, and the $m=4$ scheduler PBFT overhead. (b) Throughput: testbed achieves 5.6 tasks/sec vs.\ 6.7 tasks/sec in simulation (16\% reduction due to Raspberry Pi CPU limitations and $m=4$ scheduler overhead). (c) Detection rate: testbed maintains 97.3\% vs.\ 98.3\% in simulation (1\% degradation due to GPS noise). (d) CPU utilization: testbed Pi nodes show higher usage (22\% vs.\ 14\%) due to cryptographic overhead on ARM architecture. Error bars: 95\% confidence intervals over 50 runs.}
    \label{fig:simulation_vs_testbed}
\end{figure*}

Table~\ref{tab:testbed_summary} summarizes the comparison.

\begin{table}[h]
\centering
\caption{Testbed vs.\ Simulation: Summary of Key Metrics ($m=4$)}
\label{tab:testbed_summary}
\small
\adjustbox{width=\linewidth}{
\begin{tabular}{lccc}
\toprule
\textbf{Metric} & \textbf{Testbed} & \textbf{Simulation} & \textbf{Deviation} \\
\midrule
Consensus latency        & 175 ± 15 ms   & 152 ± 8 ms    & +15.1\% \\
Throughput               & 5.6 tasks/sec & 6.7 tasks/sec & -16.4\% \\
Byzantine detection rate & 97.3\%        & 98.3\%        & -1.0\%  \\
CPU utilization (Agent)  & 22\%          & 14\%          & +57\%   \\
\bottomrule
\end{tabular}}
\end{table}

\textbf{Key Findings.} Real-world testbed results closely match simulation predictions, with deviations $< 16\%$ across all metrics (except CPU utilization, which is hardware-specific). This validates:
\begin{itemize}
    \item \textbf{Protocol correctness}: STLC operates correctly under real network jitter, sensor noise, and hardware constraints, with $m = 4$ scheduler replicas providing full Byzantine fault tolerance ($f_s = 1$),
    \item \textbf{Simulation accuracy}: CORE emulation provides reliable predictions for large-scale scenarios ($|N| > 10$),
    \item \textbf{Practical feasibility}: Testbed latency (175 ms) meets real-time requirements ($< 500$ ms) for industrial applications.
\end{itemize}

\subsection{Ablation Study}
\label{sec:ablation}

We dissect contributions of each architectural layer by selectively disabling components ($|N|=30$, $k=3$, $m=4$, single task):

\begin{itemize}
    \item \textbf{Full STLC}: 37 messages (97.9\% reduction vs.\ PBFT's 1,770)
    \item \textbf{No spatiotemporal validation}: 37 messages (same count, but 12\% higher latency due to Byzantine retransmissions)
    \item \textbf{No task-coupling} (global domain, $k=n$): 1,770 messages (0\% reduction, equivalent to PBFT)
    \item \textbf{No adaptive weights}: 37 messages (97.9\% reduction, but 0.7\% higher latency due to unpenalized malicious nodes)
    \item \textbf{No scheduler replication} ($m=1$, single scheduler): 14 messages (99.2\% reduction, but single-point failure vulnerability and no Byzantine fault tolerance at the coordination layer)
\end{itemize}

\textbf{Key Insight}: Task-coupling is the primary driver of communication efficiency (84\% contribution), while spatiotemporal validation primarily improves latency and robustness. Scheduler replication ($m=4$ vs.\ $m=1$) adds $37 - 14 = 23$ messages (164\% overhead over $m=1$, or 2.2\% overhead relative to PBFT's 1,770) but is essential for eliminating single-point failures and providing full Byzantine fault tolerance at the coordination layer.

\textbf{Spatiotemporal Filtering Effectiveness.} Of 10,000 messages during 300 s experiment with 30\% Byzantine nodes: 1,000 (10\%) valid messages; 4,500 (45\%) timestamp violations (3,200 replay, 1,300 future-dated); 3,000 (30\%) spatial violations (2,100 kinematically infeasible, 900 stationary node movement claims); 1,500 (15\%) task inconsistency (cross-domain attacks). The 90\% filtering rate avoids $9{,}000 \times 37 = 333{,}000$ consensus messages, demonstrating that $O(1)$ validation overhead vastly outweighs the $O(k^2)$ consensus cost it prevents.

\textbf{Parameter Sensitivity.} Optimal weight parameters: $\Delta w \in [0.08, 0.15]$, $\gamma \in [0.005, 0.02]$ achieve 5--10 round isolation. Current configuration ($\Delta w = 0.1$, $\gamma = 0.01$, $k_{\text{good}} = 10$) achieves 7-round isolation (near-optimal), balancing rapid malicious detection with tolerance for transient network issues.

% ============================================================
\section{Conclusion}
\label{sec:conclusion}

This paper presents STLC (Spatiotemporal Localized Consensus with Adaptive Trust), a scalable Byzantine-resilient consensus protocol tailored for industrial multi-agent systems. STLC addresses the fundamental challenge of achieving Byzantine fault tolerance in resource-constrained, large-scale industrial IoT environments where traditional global consensus approaches are prohibitively expensive.

\subsection{Key Contributions}
\label{sec:contributions}

STLC introduces three synergistic innovations that collectively reduce communication complexity from $O(|N|^2)$ to $O(k^2)$ while maintaining strong Byzantine fault tolerance:

\textbf{(1) Spatiotemporal Security Layer.} Layer~1 achieves $O(1)$ early rejection of Byzantine messages by validating timestamps (Eq.~\ref{eq:temporal_validation}), kinematic constraints (Eq.~\ref{eq:kinematic_validation}), and task consistency (Eq.~\ref{eq:task_consistency}), filtering 65--90\% of malicious messages before they enter consensus. This lightweight validation requires no cryptographic operations, making it suitable for resource-constrained industrial devices.

\textbf{(2) Task-Coupled Consensus Domains.} Layer~2 dynamically forms localized consensus domains $C(\tau_j)$ of size $k \ll |N|$ based on task-agent proximity, reducing communication from $O(|N|^2)$ (global PBFT) to $O(k^2)$ (Theorem~\ref{thm:complexity}). For $k=3$ and $|N|=100$, this achieves 99.9\% message reduction (37 vs.\ 19,900 messages) in emulation. Experimental validation on a physical testbed ($|N|=10$) demonstrates 97.9\% reduction in practice (37 vs.\ 1,770 messages for $|N|=30$ in emulation), with 60\% latency improvement (152 ms vs.\ 382 ms) and linear throughput scaling as $|N|$ increases to 100 in emulation.

\textbf{(3) Adaptive Trust-Weighted Quorum.} Layer~3 maintains dynamic weights $w_i \in [w_{\min}, 1]$ for each node, enabling gradual isolation of Byzantine nodes within $T_{\text{isolate}} = 9$ rounds (Eq.~\ref{eq:isolation_time}) without explicit exclusion.

We prove STLC satisfies fundamental Byzantine consensus properties: \textit{safety} through weighted quorum intersection (Theorem~\ref{thm:safety}), \textit{liveness} under partial synchrony (Theorem~\ref{thm:liveness}), and $O(k^2)$ communication complexity independent of system size $|N|$ (Theorem~\ref{thm:complexity}). Physical testbed deployment on Raspberry Pi agents and Intel NUC scheduler nodes ($m = 4$ replicas, $f_s = 1$ Byzantine fault tolerance) confirms protocol correctness, achieving 175 ms consensus latency (well below the 500 ms real-time threshold) with 97.3\% Byzantine detection rate for spatiotemporal attacks. STLC demonstrates correctness and feasibility on a 10-node physical testbed, with scalability to 100 nodes validated through CORE emulation. Large-scale physical deployment ($|N| > 20$) remains a direction for future work.

\subsection{Future Work}
\label{sec:future}

\textbf{(1) Large-Scale Physical Validation.} The current testbed (6 agent nodes, 4 scheduler nodes) validates protocol correctness and functional feasibility under real hardware constraints, while scalability results for $|N| > 10$ are grounded in CORE emulation. A natural next step is hybrid deployments that combine physical hardware with emulated nodes to incrementally scale validation, and collaboration with industrial partners to construct 50--100 node testbeds in real factory environments. Such deployments would provide empirical confirmation of the emulation-based scalability predictions and quantify deployment factors (e.g., heterogeneous hardware, industrial wireless interference) not captured in CORE.

\textbf{(2) Machine Learning-Based Domain Optimization.} Reinforcement learning to dynamically adjust domain size $k$ based on real-time threat levels and task characteristics, balancing efficiency ($k=3$ for 97.9\% message reduction) against robustness ($k=7$ for $<0.1\%$ fallback trigger rate). This could also enable proactive domain pre-expansion when Byzantine activity is detected early, reducing the latency penalty of reactive expansion ($42 \pm 8$ ms per expansion round).

\textbf{(3) Formal Failure Detection Assumptions.} The current spatiotemporal validation layer assumes that correct agents report accurate positions and timestamps (within sensor noise bounds $\pm 2$ m GPS, $\pm 5$ ms clock drift). A remaining open question is \textit{completeness}: whether a Byzantine agent can evade detection by staying within kinematic bounds while still causing harm (e.g., slow drift attacks). Future work will formalize detection completeness and accuracy guarantees, and investigate whether zero-knowledge proofs for kinematic feasibility can provide cryptographic completeness without revealing exact positions---enabling deployment in privacy-sensitive industrial environments while maintaining $O(1)$ validation complexity.

\textbf{(4) Integration with Sharded and Weighted BFT Systems.} STLC's task-coupled domains share conceptual similarities with sharded BFT systems (e.g., Rapidchain, OmniLedger) and weighted Byzantine quorum systems~\cite{malkhi1998byzantine}. Future work will formally compare STLC's domain formation against shard assignment strategies, and explore combining STLC's spatiotemporal locality with DAG-based protocols (e.g., Narwhal/Tusk) to achieve $O(1)$ amortized latency within domains, potentially reducing per-task complexity from $O(k^2)$ to $O(k)$ for high-throughput scenarios.

STLC demonstrates that exploiting spatiotemporal context enables practical Byzantine consensus for industrial multi-agent systems. The current results---validated on a 10-node physical testbed and 100-node emulation---provide a foundation for further investigation into large-scale physical deployment, formal detection guarantees, and integration with complementary BFT paradigms.

\begin{thebibliography}{1}
\bibitem{bouhata2025byzantine}
D. Bouhata, H. Moumen, J. A. Mazari, and A. Bounceur, ``Byzantine fault tolerance in distributed machine learning: A survey,'' \textit{J. Exp. Theor. Artif. Intell.}, vol. 37, no. 8, pp. 1331--1389, 2025.
\bibitem{wu2024half}
H. Wu, C. Yue, Y. Fan, Y. Li, and L. Zhang, ``Half a century of distributed Byzantine fault-tolerant consensus: Design principles and evolutionary pathways,'' \textit{arXiv preprint arXiv:2407.19863}, 2024.
\bibitem{liu2024improved}
X. Liu and J. Zhu, ``An improved practical Byzantine fault tolerance algorithm for aggregating node preferences,'' \textit{Sci. Rep.}, vol. 14, no. 1, p. 31200, 2024.
\bibitem{dotan2021survey}
M. Dotan, Y.-A. Pignolet, S. Schmid, S. Tochner, and A. Zohar, ``Survey on blockchain networking: Context, state-of-the-art, challenges,'' \textit{ACM Comput. Surv.}, vol. 54, no. 5, pp. 1--34, 2021.
\bibitem{bhamare2020cybersecurity}
D. Bhamare, M. Zolanvari, A. Erbad, R. Jain, K. Khan, and N. Meskin, ``Cybersecurity for industrial control systems: A survey,'' \textit{Comput. Secur.}, vol. 89, p. 101677, 2020.
\bibitem{feldman1997optimal}
P. Feldman and S. Micali, ``An optimal probabilistic protocol for synchronous Byzantine agreement,'' \textit{SIAM J. Comput.}, vol. 26, no. 4, pp. 873--933, 1997.
\bibitem{dhatterwal2022role}
J. S. Dhatterwal, K. S. Kaswan, and R. P. Ojha, ``The role of multiagent system in industry 4.0,'' in \textit{A Roadmap for Enabling Industry 4.0 by Artificial Intelligence}. Wiley Online Library, 2022, pp. 227--246.
\bibitem{hady2025multi}
M. A. Hady, S. Hu, M. Pratama, Z. Cao, and R. Kowalczyk, ``Multi-agent reinforcement learning for resources allocation optimization: A survey,'' \textit{Artif. Intell. Rev.}, vol. 58, no. 11, p. 354, 2025.
\bibitem{vlachos2024smart}
I. Vlachos, R. M. Pascazzi, M. Ntotis, K. Spanaki, S. Despoudi, and P. Repoussis, ``Smart and flexible manufacturing systems using autonomous guided vehicles (AGVs) and the Internet of Things (IoT),'' \textit{Int. J. Prod. Res.}, vol. 62, no. 15, pp. 5574--5595, 2024.
\bibitem{margret2024blockchain}
M. K. Margret, E. G. Julie, and Y. H. Robinson, ``Blockchain-enabled resilient Byzantine fault tolerance consensus mechanism for supply chain management,'' \textit{Int. J. Web Grid Serv.}, vol. 20, no. 4, pp. 438--459, 2024.
\bibitem{li2022byzantine}
J. Li, W. Abbas, M. Shabbir, and X. Koutsoukos, ``Byzantine resilient distributed learning in multirobot systems,'' \textit{IEEE Trans. Robot.}, vol. 38, no. 6, pp. 3550--3563, 2022.
\bibitem{habibi2023performance}
T. L. Habibi and R. F. Sari, ``Performance evaluation of QUIC protocol in message replication overhead in PBFT consensus using ns-3,'' \textit{Int. J. Electr., Comput., Biomed. Eng.}, vol. 1, no. 1, pp. 44--56, 2023.
\bibitem{ye2024tp}
S. Ye, Q. Chen, F. Xiao, Z. Ke, G. Yang, and H. Ma, ``TP-BFT: A faster asynchronous BFT consensus with parallel structure,'' in \textit{Proc. Int. Conf. Algorithms Archit. Parallel Process.}, Springer, 2024, pp. 60--79.
\bibitem{arun2022scalable}
B. Arun and B. Ravindran, ``Scalable Byzantine fault tolerance via partial decentralization,'' \textit{arXiv preprint arXiv:2202.13408}, 2022.
\bibitem{wang2022bft}
X. Wang, S. Duan, J. Clavin, and H. Zhang, ``BFT in blockchains: From protocols to use cases,'' \textit{ACM Comput. Surv.}, vol. 54, no. 10s, pp. 1--37, 2022.
\bibitem{wang2024improved}
R. Wang, M. Yuan, Z. Wang, and Y. Li, ``Improved fast-response consensus algorithm based on HotStuff,'' \textit{Sensors}, vol. 24, no. 16, p. 5417, 2024.
\bibitem{wu2023learning}
D. Wu, Y. T. Xu, J. Li, M. Jenkin, E. Hossain, S. Jang, Y. Xin, C. Zhang, X. Liu, and G. Dudek, ``Learning to adapt: Communication load balancing via adaptive deep reinforcement learning,'' in \textit{Proc. IEEE Global Commun. Conf. (GLOBECOM)}, 2023, pp. 2973--2978.
\bibitem{nanapu2025d2bft}
V. Nanapu, G. Banavathu, and K. E. S. Desikan, ``D2BFT: A dual Byzantine fault tolerance approach for multi-agent drone surveillance with deep reinforcement learning,'' \textit{Comput. Netw.}, p. 111750, 2025.
\bibitem{leng2019manuchain}
J. Leng, D. Yan, Q. Liu, K. Xu, J. L. Zhao, R. Shi, L. Wei, D. Zhang, and X. Chen, ``ManuChain: Combining permissioned blockchain with a holistic optimization model as bi-level intelligence for smart manufacturing,'' \textit{IEEE Trans. Syst., Man, Cybern., Syst.}, vol. 50, no. 1, pp. 182--192, 2019.
\bibitem{pirani2022survey}
M. Pirani, A. Mitra, and S. Sundaram, ``A survey of graph-theoretic approaches for analyzing the resilience of networked control systems,'' \textit{arXiv preprint arXiv:2205.12498}, 2022.
\bibitem{karimireddy2021learning}
S. P. Karimireddy, L. He, and M. Jaggi, ``Learning from history for Byzantine robust optimization,'' in \textit{Proc. Int. Conf. Mach. Learn. (ICML)}, Virtual, 2021, pp. 5311--5319.
\bibitem{mirzadeh2022filtering}
I. Mirzadeh, M. S. Haghighi, and A. Jolfaei, ``Filtering malicious messages by trust-aware cognitive routing in vehicular ad hoc networks,'' \textit{IEEE Trans. Intell. Transp. Syst.}, vol. 24, no. 1, pp. 1134--1143, 2022.
\bibitem{danezis2022narwhal}
G. Danezis, L. Kokoris-Kogias, A. Sonnino, and A. Spiegelman, ``Narwhal and tusk: A DAG-based mempool and efficient BFT consensus,'' in \textit{Proc. 17th Eur. Conf. Comput. Syst.}, 2022, pp. 34--50.
\bibitem{spiegelman2022bullshark}
A. Spiegelman, N. Giridharan, A. Sonnino, and L. Kokoris-Kogias, ``Bullshark: The partially synchronous version,'' \textit{arXiv preprint arXiv:2209.05633}, 2022.
\bibitem{gelashvili2022jolteon}
R. Gelashvili, L. Kokoris-Kogias, A. Sonnino, A. Spiegelman, and Z. Xiang, ``Jolteon and ditto: Network-adaptive efficient consensus with asynchronous fallback,'' in \textit{Proc. Int. Conf. Financial Cryptogr. Data Secur.}, Cham, Switzerland: Springer, 2022, pp. 296--315.
\bibitem{malkhi1998byzantine}
D. Malkhi and M. Reiter, ``Byzantine quorum systems,'' \textit{Distributed Computing}, vol. 11, no. 4, pp. 203--213, 1998.
\bibitem{mivsic2020adapting}
J. Mi\v{s}i\'c, V. B. Mi\v{s}i\'c, X. Chang, and H. Qushtom, ``Adapting PBFT for use with blockchain-enabled IoT systems,'' \textit{IEEE Trans. Veh. Technol.}, vol. 70, no. 1, pp. 33--48, 2020.
\bibitem{iec61508}
International Electrotechnical Commission, ``Functional safety of electrical/electronic/programmable electronic safety-related systems - Part 1: General requirements,'' IEC Standard 61508-1:2010, 2010.
\bibitem{bessani2014state}
A. Bessani, J. Sousa, and E. E. P. Alchieri, ``State machine replication for the masses with BFT-SMART,'' in \textit{Proc. 44th Annu. IEEE/IFIP Int. Conf. Dependable Syst. Netw.}, 2014, pp. 355--362.
\bibitem{tsaliki2024accelerating}
S. R. Tsaliki, ``Accelerating ETCD insertion operations through advanced complexity reduction,'' \textit{arXiv preprint}, 2024.
\bibitem{carr2022towards}
H. Carr, C. Jenkins, M. Moir, V. C. Miraldo, and L. Silva, ``Towards formal verification of HotStuff-based Byzantine fault tolerant consensus in Agda,'' in \textit{Proc. NASA Formal Methods Symp.}, Cham, Switzerland: Springer, 2022, pp. 616--635.
\end{thebibliography}

\end{document}